%%%%%%%%%%%%%%%%%%%%%%%%%%%%%%%%%%%%%%%%%
% Beamer Presentation
% LaTeX Template
% Version 1.0 (10/11/12)
%
% This template has been downloaded from:
% http://www.LaTeXTemplates.com
%
% License:
% CC BY-NC-SA 3.0 (http://creativecommons.org/licenses/by-nc-sa/3.0/)
%
%%%%%%%%%%%%%%%%%%%%%%%%%%%%%%%%%%%%%%%%%

%----------------------------------------------------------------------------------------
%	PACKAGES AND THEMES
%----------------------------------------------------------------------------------------

\documentclass[10pt]{beamer}

\mode<presentation> {

% The Beamer class comes with a number of default slide themes
% which change the colors and layouts of slides. Below this is a list
% of all the themes, uncomment each in turn to see what they look like.

%\usetheme{default}
%\usetheme{AnnArbor}
\usetheme{Antibes}  %Good
%\usetheme{Bergen}
%\usetheme{Berkeley}
%\usetheme{Berlin}
%\usetheme{Boadilla}  %Good
%\usetheme{CambridgeUS}
%\usetheme{Copenhagen}
%\usetheme{Darmstadt}
%\usetheme{Dresden}
%\usetheme{Frankfurt}
%\usetheme{Goettingen}
%\usetheme{Hannover}
%\usetheme{Ilmenau}
%\usetheme{JuanLesPins}
%\usetheme{Luebeck}
%\usetheme{Madrid}  %Good
%\usetheme{Malmoe}
%\usetheme{Marburg}
%\usetheme{Montpellier}
%\usetheme{PaloAlto}
%\usetheme{Pittsburgh}
%\usetheme{Rochester}  %Good
%\usetheme{Singapore}  %Good
%\usetheme{Szeged}
%\usetheme{Warsaw}

% As well as themes, the Beamer class has a number of color themes
% for any slide theme. Uncomment each of these in turn to see how it
% changes the colors of your current slide theme.

%\usecolortheme{albatross}
%\usecolortheme{beaver}
%\usecolortheme{beetle}
%\usecolortheme{crane}
\usecolortheme{dolphin}  %Good
%\usecolortheme{dove}
%\usecolortheme{fly}
%\usecolortheme{lily}
%\usecolortheme{orchid}  %Good
%\usecolortheme{rose}  %Good
%\usecolortheme{seagull}
%\usecolortheme{seahorse}  %Good
%\usecolortheme{whale}  %Good
%\usecolortheme{wolverine}

%\setbeamertemplate{footline} % To remove the footer line in all slides uncomment this line
%\setbeamertemplate{footline}[page number] % To replace the footer line in all slides with a simple slide count uncomment this line

%\setbeamertemplate{navigation symbols}{} % To remove the navigation symbols from the bottom of all slides uncomment this line
}

\usepackage{graphicx} % Allows including images
\usepackage{booktabs} % Allows the use of \toprule, \midrule and \bottomrule in tables
\usepackage{amsmath}
\usepackage{listings}
\usepackage{courier}           % for the courier font (optional)
\usepackage{multicol}          % for two equations side by side
\usepackage{centernot}
\usepackage{hyperref}
\usepackage{xcolor}
\usepackage{stmaryrd}
\usepackage{hieroglf}
\usepackage{scalerel}
\usepackage{enumerate}
\usepackage{courier}
\usepackage{appendixnumberbeamer}
\usepackage{pgfplots}
\usepackage{mdframed} %nice frames
\usepackage{bm}

\hypersetup{ % play with these to change the look of hyperlinks
	colorlinks=true,
%	linkcolor=white,
%	filecolor=white,
  allcolors=.,
	urlcolor=blue,
%	citecolor=black
}


%\setbeamertemplate{bibliography item}{\insertbiblabel}


%------------------------------------------------------------------------------------------
%    ENVIRONMENT SETTINGS
%------------------------------------------------------------------------------------------

% TEMPLATE SETTINGS
\setbeamertemplate{theorems}[numbered]
\setbeamertemplate{sidebar right}{}
\setbeamertemplate{footline}{%
	\hspace{1em}\small\insertframenumber/\inserttotalframenumber
	\hfill\usebeamertemplate***{navigation symbols}}

\makeatletter
\renewcommand{\itemize}[1][]{%
  \beamer@ifempty{#1}{}{\def\beamer@defaultospec{#1}}%
  \ifnum \@itemdepth >2\relax\@toodeep\else
    \advance\@itemdepth\@ne
    \beamer@computepref\@itemdepth% sets \beameritemnestingprefix
    \usebeamerfont{itemize/enumerate \beameritemnestingprefix body}%
    \usebeamercolor[fg]{itemize/enumerate \beameritemnestingprefix body}%
    \usebeamertemplate{itemize/enumerate \beameritemnestingprefix body begin}%
    \list
      {\usebeamertemplate{itemize \beameritemnestingprefix item}}
      {\setlength\itemsep{1em}% NEW
        \def\makelabel##1{%
          {%
            \hss\llap{{%
                \usebeamerfont*{itemize \beameritemnestingprefix item}%
                \usebeamercolor[fg]{itemize \beameritemnestingprefix item}##1}}%
          }%
        }%
      }
  \fi%
  \beamer@cramped%
  \raggedright%
  %\beamer@firstlineitemizeunskip%
}
\makeatother

% ENUMERATE COLOR SETTINGS
\newenvironment{doneenv}{\only{\setbeamercolor{local structure}{fg=gray}\color{gray}}}{}
\newenvironment{comeenv}{\only{\setbeamercolor{local structure}{fg=red}\color{red}}}{}

% LSTLISTING SETTINGS
\lstdefinestyle{myStyle}{
	%	language=Coq,
	keywords={Inductive,Require,Import,Definition,Fixpoint,match,with,end,let,in,fix},
	basicstyle=\footnotesize\ttfamily,
	keywordstyle=\color{green}, % Blue clashes with the cyan links. Change if you want.
	stepnumber=1,
	tabsize=2,
	numbers=none,
	numberstyle=\tiny,
	numbersep=5pt,
	showspaces=false,
	escapechar=`,
	showstringspaces=false,
	backgroundcolor = \color{lightgray},
}

\lstset{basicstyle=\ttfamily\footnotesize,breaklines=true}
%basicstyle=\fontsize{10}{11}\selectfont\ttfamily,

\lstdefinestyle{myTinyStyle}{
	%   language=Coq,
	keywords={Inductive,Require,Import,Definition,Fixpoint,match,with,end,let,in,fix},
	basicstyle=\normalfont\fontsize{8}{8.3}\tt,
	keywordstyle=\color{green}, % Blue clashes with the cyan links. Change if you want.
	stepnumber=1,
	tabsize=2,
	numbers=none,
	numberstyle=\tiny,
	numbersep=5pt,
	showspaces=false,
	escapechar=`,
	showstringspaces=false
}

\lstset{style=myStyle}
\makeatletter
\newlength{\@mli}
\newcommand{\mli}[1]{%
	\settowidth{\@mli}{\lstinline/#1/}
	\hspace{-.5ex}\begin{minipage}[t]{\@mli}\lstinline/#1/\end{minipage}}
\makeatother
\newcommand{\li}[1]{\ifmmode\mbox{\mli{#1}}\else\mbox{\lstinline/#1/}\fi}

\newcommand{\coq}{\scalebox{.6}{\textpmhg{\Ha}}}

\newcommand{\cdt}[2]{#1\mathrel{\mathchoice{\mkern-5mu}{\mkern-5mu}{\mkern-2mu}{\mkern-2mu}} ^{\mathcal{C}}_{#2} \mathrel{\mathchoice{\mkern-3.5mu}{\mkern-3.5mu}{\mkern-2mu}{\mkern-2mu}}}

\newcommand{\cdw}[2]{#1\mathrel{\mathchoice{\mkern-5mu}{\mkern-5mu}{\mkern-2mu}{\mkern-2mu}} ^{\mathcal{CW}}_{#2} \mathrel{\mathchoice{\mkern-3.5mu}{\mkern-3.5mu}{\mkern-2mu}{\mkern-2mu}}}

\newcommand{\cds}[2]{#1\mathrel{\mathchoice{\mkern-5mu}{\mkern-5mu}{\mkern-2mu}{\mkern-2mu}} ^{\mathcal{CS}}_{#2} \mathrel{\mathchoice{\mkern-3.5mu}{\mkern-3.5mu}{\mkern-2mu}{\mkern-2mu}}}

\newcommand{\rf}[2]{{#1\mathrel{\mathchoice{\mkern-5mu}{\mkern-5mu}{\mkern-2mu}{\mkern-2mu}} ^{\mathcal{R}}_{#2} \mathrel{\mathchoice{\mkern-3.5mu}{\mkern-3.5mu}{\mkern-2mu}{\mkern-2mu}}}}

\newcommand{\invhyper}{I}
\renewcommand{\angle}[1]{\langle #1\rangle}

\newcommand{\W}{\mathcal{W}}
\newcommand{\contract}{\text{\scshape{cont}}}
\newcommand{\repeatable}{\text{\scshape{rept}}}

\newcommand{\divc}{\text{divc}}
\newcommand{\logc}{\text{logc}}
\newcommand{\logstarc}{\text{logc*}}

\newcommand{\runtime}{\mathcal{T}}
\newcommand{\Texp}{\runtime_{\li{exp}}}
\newcommand{\Tmul}{\runtime_{\li{mul}}}
\newcommand{\Tleb}{\runtime_{\li{leb}}}
\newcommand{\Tsucc}{\runtime_{\li{succ}}}
\newcommand{\bigO}{\text{O}}

\newcommand{\Ack}{\ensuremath{\mathcal{A}}}
\newcommand{\CA}{\ensuremath{\text{C}}}
%\renewcommand{\S}{\ensuremath{\text{S}}}
\newcommand{\CR}{\ensuremath{\text{CR}}}
\newcommand{\CRH}{\ensuremath{\text{CRH}}}
\newcommand{\wt}[1]{\ensuremath{\widetilde{#1}}}
\newcommand{\IAR}{\ensuremath{\text{IAR}}}
\newcommand{\IARH}{\ensuremath{\text{IARH}}}

\newcommand\hide[1]{}
\newenvironment{centermath}
{\begin{center}$\displaystyle}
	{$\end{center}}
\renewcommand{\note}[2][polish]{{\color{red} #2}{\marginpar{\tiny \color{blue} #1}}}
\renewcommand{\implies}{\Rightarrow}
\renewcommand{\iff}{\Leftrightarrow}


% TABLE, ARRAY SETTINGS
\newcommand{\topspace}[1]{\rule{0pt}{2.6ex+#1}}        % = `top' strut
\newcommand{\botspace}[1]{\rule[-0.9ex-#1]{0pt}{0pt}}   % = `bottom' strut


% DARKENED COLORS
\colorlet{red}{red!80!black}
\colorlet{blue}{blue!80!black}
\colorlet{green}{green!50!black}
\colorlet{lightgray}{gray!30!white}

% FBOX SETTINGS
\newcommand{\imppar}[1]{{\setlength{\fboxrule}{1pt}\fcolorbox{red}{white}{#1}}}
\newcommand{\impinline}[1]{{\setlength{\fboxrule}{1pt}\setlength{\fboxsep}{2pt}\fcolorbox{white}{lightgray}{#1}}}

%\AtBeginSection[]{
%	\begin{frame}
%	\vfill
%	\centering
%	\begin{beamercolorbox}[sep=8pt,center,shadow=true,rounded=true]{title}
%		\Huge{\insertsectionhead}\par%
%	\end{beamercolorbox}
%	\vfill
%	\end{frame}
%}

%\AtBeginSection[*]{
%	\begin{frame}
%	\vfill
%	\centering
%	\begin{beamercolorbox}[sep=8pt,center,shadow=true,rounded=true]{title}
%		\Huge{\insertsectionhead}\par%
%	\end{beamercolorbox}
%	\vfill
%\end{frame}
%}
%----------------------------------------------------------------------------------------
%	TITLE PAGE
%----------------------------------------------------------------------------------------

\title[Inverting the Ackermann Hierarchy]{A Functional Proof Pearl: \protect\\ Inverting the Ackermann Hierarchy} % The short title appears at the bottom of every slide, the full title is only on the title page
\author{Linh Tran, Anshuman Mohan, \underline{Aquinas Hobor}} % Your name
\institute[NUS]{$\includegraphics[scale=0.45]{NUS_logo_full-horizontal.jpg}$}
%{$\includegraphics[scale=0.33]{NUS_logo_full-horizontal.jpg}} % Your institution as it will appear on the bottom of every slide, may be shorthand to save space
% {
% National University of Singapore % Your institution for the title page
% }
\date{University College London\\ \today} % Date, can be changed to a custom date

\begin{document}

\begin{frame}[plain]
\titlepage % Print the title page as the first slide
\end{frame}

%----------------------------------------------------------------------------------------
%	PRESENTATION SLIDES
%----------------------------------------------------------------------------------------
% ============================================================================
% OVERVIEW
% ============================================================================

\section*{Introduction}
\label{sec: overview}
\begin{frame}
\frametitle{The Ackermann Function}
	
	The Ackermann-P\'eter function is defined as:
	\begin{equation*}
	A(n, m) \triangleq \begin{cases}
	m + 1 & \text{ when } n = 0 \\
	A(n-1, 1) & \text{ when } n > 0, m = 0 \\
	A\big(n-1, A(n, m-1)\big) & \text{ otherwise}
	\end{cases}
	\end{equation*}
	
	\pause 
	The \emph{diagonal} Ackermann function is $\Ack(n)~\triangleq~A(n, n)$.
	
	\bigskip
	
	\pause 
	First few values of $\Ack(n)$:
	
  \begin{minipage}{0.5\linewidth}
		\begin{equation*}
	\begin{array}{r|ccccc}
	 n & 0 & 1 & 2 & 3 & 4 \\ \hline
	 \Ack(n) & 1 & 3 & 7 & 61 & 2^{2^{2^{65536}}} - 3 \topspace{3pt}
	\end{array}
	\end{equation*}
  \end{minipage}
  \quad \pause 
  \begin{minipage}{0.4\linewidth}
  	\textcolor{red}{Explosive growth!}
  \end{minipage}

\end{frame}


%\begin{frame}
%\frametitle{Initial values for $\Ack(n)$ and $\alpha(n)$}
%\begin{columns}[T]
%	\begin{column}{0.4\textwidth}
%		
%		\begin{equation*}
%		\begin{array}{|ll|}
%		n & \Ack(n) \\
%		0 & 1 \\
%		1 & 3 \\
%		2 & 7 \\
%		3 & 61 \\
%		4 & 2^{2^{2^{65536}}} - 3
%		\end{array}
%		\end{equation*}
%		
%		Grows astronomically fast!
%	\end{column}
%
%  \begin{column}{0.5\textwidth}
%  	
%  	\begin{equation*}
%  	\begin{array}{|ll|}
%  	n & \alpha(n) \\
%  	0, 1 & 0 \\
%  	2, 3 & 1 \\
%  	4, 5, 6, 7 & 2 \\
%  	8, 9, \ldots, 61 & 3 \\
%  	62, \ldots, 2^{2^{2^{65536}}} - 3 & 4
%  	\end{array}
%  	\end{equation*}
%  	
%  	Grows astronomically slowly!
%  \end{column}
%\end{columns}
%\end{frame}


% \begin{frame}
% \frametitle{Introduction: Growth Patterns}

%  First few values of $\Ack(n)$ and $\alpha(n)$:
		
% %		\begin{equation*}
% %		\begin{array}{|ll|c|ll|}
% %		n & \Ack(n) & \hspace{3em} & n & \alpha(n) \\[3pt]
% %		0 & 1 & & 0, 1 & 0 \\[3pt]
% %		1 & 3 & & 2, 3 & 1 \\[3pt]
% %		2 & 7 & & 4, 5, 6, 7 & 2 \\[3pt]
% %		3 & 61 & & 8, 9, \ldots, 61 & 3 \\[3pt]
% %		4 & 2^{2^{2^{65536}}} - 3 & & 62, 63 \ldots, 2^{2^{2^{65536}}} - 3 & 4 \\[5pt]
% %		\multicolumn{2}{|l|}{\text{Grows astronomically fast!}} & & \multicolumn{2}{|l|}{\text{Grows astronomically slow!}}
% %		\end{array}
% %		\end{equation*}
		
% 		\begin{tikzpicture}
% 		\begin{axis}[
% %		symbolic x coords={a small bar,a medium bar,a large bar},
%     yticklabels={,,},
% 		xtick=data]
% 		\addplot coordinates {
% 			(0, 1)
% 			(1, 3)
% 			(2, 7)
% 			(3, 61)
% 			(4, 9999999)
			
% 		};
% 		\end{axis}
% 		\end{tikzpicture}

% %TODO Fix this graph? Graph for alpha? Figure with caption?

% \end{frame}


\begin{frame}
\frametitle{The Inverse Ackermann Function}

The \emph{inverse Ackermann function}, $\alpha$, maps $n$ to the smallest~$k$ for
which~$n \le \Ack(k)$, \emph{i.e} \impinline{$\alpha(n) \triangleq \min\left\{k\in \mathbb{N} : n \le \Ack(k)\right\}$}.

\smallskip

\pause 
$\alpha(n)$ grows slowly but is hard to \emph{compute} for large $n$
\\ because it is entangled with the explosively-growing $\Ack(k)$.

\bigskip


%\textbf{Naive Approach:} starting at $k=0$, calculate $\Ack(k)$,
%compare it to $n$, \\ and increment $k$ until $n \le \Ack(k)$.
\pause 
\textbf{Naive Approach:} Compute $\Ack(0), \Ack(1), \ldots$ until $n\le \Ack(k)$. Return $k$.

\bigskip

\pause 
\textbf{Time complexity:} $\Omega(\Ack(\alpha(n)))$.
\\ Computing $\alpha(100) \mapsto^{*} 4$ requires at least
%$\Ack(4) = 2^{2^{2^{65536}}}$ steps!
$\Ack(4) = 2^{2^{2^{65536}}} - 3$ steps!

\bigskip

\pause 
\textbf{Engineering hack:} Hardcode with lookup tables. $n > 61 \implies \text{ans} = 4$.

\bigskip

\pause 
\imppar{\text{\textbf{Our Goal.} Compute $\alpha$ for \emph{all inputs} without computing $\Ack$.}}
\end{frame}

%\begin{frame}
%\frametitle{Introduction: Ackermann \emph{vs} Hyperoperations}
%The Ackermann function is easy to define, but hard to
%understand.
%
%We see it as
%a sequence of $n$-indexed functions $\Ack_n \triangleq \lambda b.A(n,b)$, where for each $n>0$, $\Ack_n$ is the result of applying the previous $\Ack_{n-1}$ $b$ times.
%
%%with a
%%\href{https://github.com/inv-ack/inv-ack/blob/7270e64a2600b771f2b1b1b151f7d13fb2ae6c97/repeater.v\#L161-L177}{\emph{kludge}}. %Linked by Linh
%
%\bigskip
%
%The hierarchical structure resembles that of \emph{hyperoperations}.
%
%%To better understand the Ackermann function as a hierarchy and this kludge, we explore the closely-related hyperoperations.
%
%\end{frame}


\begin{frame}[fragile]
\frametitle{Our Solution}

\vspace{-1em}
\lstset{style=myTinyStyle}
% Linked by A
\begin{mdframed}[backgroundcolor=lightgray, roundcorner=10pt,leftmargin=0, rightmargin=0, innerleftmargin=0, innertopmargin=-5,innerbottommargin=-5, outerlinewidth=0, linecolor=lightgray]
\begin{lstlisting}
Require Import Omega Program.Basics.

`\href{https://github.com/inv-ack/inv-ack/blob/7270e64a2600b771f2b1b1b151f7d13fb2ae6c97/inv_ack_standalone.v#L6-L11}{Fixpoint cdn\_wkr}` (a : nat) (f : nat -> nat) (n b : nat) :=
 match b with 0 => 0 | S b' =>
  if (n <=? a) then 0 else S (cdn_wkr f a (f n) k')
 end.

`\href{https://github.com/inv-ack/inv-ack/blob/7270e64a2600b771f2b1b1b151f7d13fb2ae6c97/inv_ack_standalone.v#L14}{Definition countdown\_to}` a f n := cdn_wkr a f n n.

`\href{https://github.com/inv-ack/inv-ack/blob/7270e64a2600b771f2b1b1b151f7d13fb2ae6c97/inv_ack_standalone.v#L32-L38} {Fixpoint inv\_ack\_wkr}` (f : nat -> nat) (n k b : nat) :=
 match b with 0 => 0 | S b' =>
  if (n <=? k) then k else let g := (countdown_to f 1) in
                      inv_ack_wkr (compose g f) (g n) (S k) b
 end.

`\href{https://github.com/inv-ack/inv-ack/blob/7270e64a2600b771f2b1b1b151f7d13fb2ae6c97/inv_ack_standalone.v#L42-L46}{Definition inv\_ack\_linear}` (n : nat) : nat :=
 match n with 0 | 1 => 0 | _ => 
  let f := (fun x => x - 2) in inv_ack_wkr f (f n) 1 (n - 1)
 end.
\end{lstlisting}
\end{mdframed} 
\end{frame}


\begin{frame}
\frametitle{Ackermann \emph{vs} Hyperoperation}

The Ackermann function is easy to define, but hard to
understand.

\bigskip

\pause 	
Let's index by the first argument. \\ \smallskip
Define $\Ack_n \triangleq \lambda b.A(n,b)$. \\ \smallskip
Then, for $n>0$, $\Ack_n$ is the result of applying the previous $\Ack_{n-1}$ $b$ times.

%with a
%\href{https://github.com/inv-ack/inv-ack/blob/7270e64a2600b771f2b1b1b151f7d13fb2ae6c97/repeater.v\#L161-L177}{\emph{kludge}}. %Linked by Linh

\bigskip

\pause 
The hierarchical structure resembles that of \textcolor{red}{\emph{hyperoperations}}.

\smallskip
Studying hyperoperations helps us understand the Ackermann hierarchy.

%To better understand the Ackermann function as a hierarchy and this kludge, we explore the closely-related hyperoperations.
\end{frame}


\begin{frame}
\frametitle{Introduction: Ackermann \emph{vs} Hyperoperation}

Treating $b$ as the main argument, we can build their \emph{upper inverses}:

\begin{table}[t]
	\begin{centermath}
		\begin{array}{c@{\hskip 0.5em}|@{\hskip 1em}c@{\hskip 1em}c@{\hskip 1em}|@{\hskip 1em}c@{\hskip 1em}c}
%			  & \multicolumn{2}{|@{\hskip 0.5em}c@{\hskip 0.5em}|}{\text{Main hierarchies}} & \multicolumn{2}{|@{\hskip 0.5em}c@{\hskip 0.5em}|}{\text{Inverses hierarchies}} \\
			n & a [n] b & \Ack_n(b) & a \angle{n} b & \alpha_n(b)\\
			\hline
			0 & 1 + b & 1 + b & b - 1 & b - 1 \\
			1 & a + b & 2 + b & b - a & b - 2 \\
			2 & a \cdot b & 2b + 3 & \left\lceil \frac{b}{a} \right\rceil & \left\lceil \frac{b-3}{2} \right\rceil \\
			3 & a^b & 2^{b + 3} - 3 & \left\lceil \log_a ~ b \right\rceil & \left\lceil \log_2 ~ (b + 3)\right\rceil - 3 \\
			[1pt]
			4 & \underbrace{a^{.^{.^{.^a}}}}_b & \underbrace{2^{.^{.^{.^2}}}}_{b+3} - 3 & \log^*_a ~ b & \log^*_2 ~ (b + 3) - 3
		\end{array}
	\end{centermath}
	\label{table: hyperop-ack-inv}
\end{table}

\pause
Connection?
\pause
\\ \href{https://github.com/inv-ack/inv-ack/blob/7270e64a2600b771f2b1b1b151f7d13fb2ae6c97/repeater.v\#L161-L177}{Aha!} 
$\Ack_n(b) = {\color{red}2}[n](b{\color{red}+3}) {\color{red}- 3}$
\pause
\quad and \quad $\alpha_n(b) = {\color{red}2}\angle{n}(b{\color{red}+3}) {\color{red}- 3}$.

\end{frame}


\begin{frame}
\frametitle{Introduction: Inverse Hierarchies to Inverse Ackermann}

We explore the upper inverse relation:
\begin{equation*}
\begin{cases}
\forall b. \forall c.\quad b \le \Ack_n(c) & \!\! \iff \ \ \alpha_n(b)\le c \\
\forall b. \forall c. \quad b \le a[n]c & \!\! \iff \ \ a\angle{n}b \le c
\end{cases}
\end{equation*}

\textbf{Redefine $\bm{\alpha}$:}
$\alpha(n) = \min\{k: n\le \Ack_k(k) \} = \min\{k: \alpha_k(n)\le k \}$.

\bigskip

\pause
\textbf{Computing $\bm{\alpha} $ through $\bm{\alpha_i}$!} No need to go through $\Ack$.

\bigskip

\pause
\imppar{\textbf{Goal.} Build the inverse towers independent from the original towers.}

\end{frame}






%\begin{frame}
%\frametitle{}
%\end{frame}

%TODO: think of a better way to make the roadmap i.e. diagram or something

\section*{Roadmap}
\begin{frame}
\frametitle{Roadmap}
\setlength{\leftmargini}{3.8em}
\setbeamertemplate{enumerate items}[default]
\setbeamercolor{local structure}{fg=black}

\imppar{\textbf{Goal.} Inverting $\Ack$ - without computing $\Ack$.}

\bigskip

\begin{enumerate}[\bfseries Step 1.]
  \itemsep 2ex
  \pause
	\item<come@6>
	Build the hyperoperations/Ackermann hierarchies via \textbf{Repeater}.
	
	\item<3->
	Build inverses of the hyperoperations/Ackermann hierarchies \\ via \textbf{Countdown}, the inverse of \emph{Repeater}.
	
	\item<4->
	Use the hierarchy to implement the inverse Ackermann function \\ for inputs encoded in unary. $\bigO (n)$.
	
	\item<5-> [\textbf{Bonus.}]
	Improve efficiency for inputs encoded in binary. $\bigO (\log_2n)$.	
\end{enumerate}
\end{frame}


\section{Step 1: Hyperoperations and Ackermann via Repeater}
\label{sec: countdown-repeater}
Let us formally define hyperoperations and clarify the intuition
given by Table~\ref{table: hyperop-ack-inv} by relating hyperoperations to
the Ackermann function.
The first hyperoperation (level 0) is simply successor, and
every hyperoperation that follows is the repeated application of the previous.
Level 1 is thus addition, and $b$ repetitions of addition
give level 2, multiplication. Next, $b$ repetitions of
multiplication give level 3, exponentiation.
There is a subtlety here: in the former case, we add $a$
repeatedly to the additive identity $0$, but in the
latter case, we multiply $a$ repeatedly to the multiplicative identity $1$. 
The formal definition of hyperoperation is:
%\marginpar{\tiny \color{blue} Should 2 just be under 1?}
%\begin{equation}
%\label{eq:hyper}
%\begin{array}{lrcl}
%\textit{1. 0$^{\textit{th}}$ level: } & a[0]b & ~ \triangleq ~ & b + 1 \\
%\textit{2. Initial values: } & a[n+1]0 & ~ \triangleq ~ &
%\begin{cases}
%a & \text{when } n = 0 \\
%0 & \text{when } n = 1 \\
%1 & \text{otherwise}
%\end{cases} \\
%\textit{3. Recursive rule: } \quad & a[n+1](b+1) & ~ \triangleq ~ & a[n]\big(a[n+1]b\big)
%\end{array}
%\end{equation}
% Edited by Linh
\begin{equation}
\label{eq:hyper}
\begin{array}{lrcl}
\textit{1. 0$^{\textit{th}}$ level:} & a[0]b & \triangleq & b + 1 \\
\textit{2. Initial values:} & a[n+1]0 & \triangleq &
  \begin{cases}
    a & \text{when } n = 0 \\
    0 & \text{when } n = 1 \\
    1 & \text{otherwise}
  \end{cases} \\
\textit{3. Recursive rule:} & a[n+1](b+1) & \triangleq & a[n]\big(a[n+1]b\big)
\end{array}
\end{equation}
The recursive rule looks complicated, but is actually just \emph{repeated application} in disguise. By fixing~$a$ and treating~$a[n]b$ as a function of~$b$, we can write
% \begin{equation*}
% \begin{array}{lll}
% a[n+1]b & ~ = ~ a[n]\big(a[n+1](b-1)\big) & ~ = ~ a[n]\big(a[n](a[n+1](b-2))\big) \\
%  & ~ = ~ \underbrace{\big( a[n]\circ a[n]\circ \cdots \circ a[n] \big)}_{b \text{ times}} \big(a[n+1]0\big) & ~ = ~ \big(a[n]\big)^{(b)}\big(a[n+1]0\big)
% \end{array}
% \end{equation*}
\begin{equation*}
\begin{array}{lll}
a[n+1]b ~& = ~ a[n]\big(a[n+1](b-1)\big) \\
         & = ~ a[n]\big(a[n](a[n+1](b-2))\big) \\
         & = ~ \underbrace{\big( a[n]\circ a[n]\circ \cdots \circ a[n] \big)}_{b \text{ times}} \big(a[n+1]0\big) \\
         & = ~ \big(a[n]\big)^{(b)}\big(a[n+1]0\big)
\end{array}
\end{equation*}
%where $f^{(k)}(u) ~ \triangleq ~ \overbrace{(f\circ f\circ \cdots \circ f)}^{k \text{ times}} (u)$,
where $f^{(k)}(u) \triangleq ~ (f\circ f\circ \cdots \circ f)(u)$ denotes $k$ compositional applications of a function~$f$ to an
input~$u$; $f^{(0)}(u) = u$ \linebreak (\emph{i.e.} applying $0$ times yields the identity).

This insight helps us encode hyperoperations~(\ref{eq:hyper}) and
Ackermann~(\ref{eq:ackermann}) in Coq.  
Notice that the recursive case of hyperoperations and
the third case of Ackermann both feature deep nested recursion, 
which makes our task tricky. 
In the outer recursive call, the first argument is shrinking
but the second is expanding explosively; in the inner recursive call, the first argument is
constant but the second is shrinking. The elegant solution uses double recursion~\cite{bertotcast} as follows:
% Linked by A
\begin{lstlisting}
`\href{https://github.com/inv-ack/inv-ack/blob/7270e64a2600b771f2b1b1b151f7d13fb2ae6c97/repeater.v#L51-L52}{\color{blue}Definition hyperop\_init}` (a n : nat) : nat :=
  match n with 0 => a | 1 => 0 | _ => 1 end.

`\href{https://github.com/inv-ack/inv-ack/blob/7270e64a2600b771f2b1b1b151f7d13fb2ae6c97/repeater.v#L55-L64}{\color{blue}Fixpoint hyperop\_original}` (a n b : nat) : nat :=
  match n with
  | 0    => 1 + b
  | S n' => let fix hyperop' (b : nat) :=
             match b with
             | 0    => hyperop_init a n'
             | S b' => hyperop_original a n' 
                                  (hyperop' b')
              end in hyperop' b
  end.

`\href{https://github.com/inv-ack/inv-ack/blob/7270e64a2600b771f2b1b1b151f7d13fb2ae6c97/repeater.v#L123-L132}{\color{blue}Fixpoint ackermann\_original}` (m n : nat) : nat :=
  match m with
  | 0    => 1 + n
  | S m' => let fix ackermann' (n : nat) : nat :=
             match n with
             | 0    => ackermann_original m' 1
             | S n' => ackermann_original m' 
                                 (ackermann' n')
             end in ackermann' n
  end.
\end{lstlisting}
Coq is satisfied since both recursive calls are on structurally smaller arguments.
Moreover, our encoding makes the structural similarities
 readily apparent.  In fact, the only essential difference is the initial values
(\emph{i.e.} the second case of both definitions): the Ackermann function uses $\Ack(n-1,1)$, whereas
hyperoperations use the initial values given in~\eqref{eq:hyper}.

We notice that the deep recursion in both cases is expressing the same notion
of repeated application, and this leads us to another useful idea. We can elegantly express the relationship
between the $(n+1)^{\text{th}}$ and $n^{\text{th}}$ levels via a higher-order function that transforms the latter level
to the former using a version of the well-known function iterator
\li{iter}~\cite{bertotcast}:
\begin{defn}
$\forall a\in \mathbb{N}, f: \mathbb{N}\to \mathbb{N}$, the
\href{https://github.com/inv-ack/inv-ack/blob/7270e64a2600b771f2b1b1b151f7d13fb2ae6c97/repeater.v#L32-L36}{\color{blue}\emph{repeater from}}
$a$ of $f$, denoted by $\rf{f}{a}$ , is a function $\mathbb{N}\to \mathbb{N}$ such that $\rf{f}{a}(n) = f^{(n)}(a)$.
% Linked by A
%\begin{lstlisting}
%`\href{https://github.com/inv-ack/inv-ack/blob/7270e64a2600b771f2b1b1b151f7d13fb2ae6c97/repeater.v#L32-L36}{Fixpoint repeater\_from}` (f : nat -> nat) (a n : nat) : nat :=
%  match n with 0 => a | S n' => f (repeater_from f a n') end.
%\end{lstlisting}
\begin{lstlisting}
`\href{https://github.com/inv-ack/inv-ack/blob/7270e64a2600b771f2b1b1b151f7d13fb2ae6c97/repeater.v#L32-L36}{\color{blue}Fixpoint repeater\_from}` f a n :=
  match n with
  | 0 => a
  | S n' => f (repeater_from f a n') 
  end.
\end{lstlisting}
\end{defn}
\noindent This allows simple and function-oriented definitions of hyperoperations and the
Ackermann function that we give below. Note that the Curried $a[n-1]$ denotes
the single-variable function $\lambda b.a[n-1]b$.
%\begin{equation*}
%\vspace{-0.75em}
%a[n]b ~ \triangleq ~ \begin{cases}
%b + 1 & \text{when } n = 0 \\
%\rf{a[n-1]}{a_{n-1}}(b) & \text{otherwise}
%\end{cases}
%\quad \text{ where } \ a_n ~ \triangleq ~ \begin{cases}
%a & \text{when } n = 0 \\
%0 & \text{when } n = 1 \\
%1 & \text{otherwise}
%\end{cases}
%\end{equation*}
% Edited by Linh
\begin{equation*}
a[n]b ~ \triangleq ~ \begin{cases}
b + 1 & \text{when } n = 0 \\
\rf{a[n-1]}{a_{n-1}}(b) & \text{otherwise}
\end{cases}
\end{equation*}
\begin{equation*}
\hspace{8em}\text{where  } a_n ~ \triangleq ~ \begin{cases}
a & \text{when } n = 0 \\
0 & \text{when } n = 1 \\
1 & \text{otherwise}
\end{cases}
\end{equation*}
% Linked by Anshuman
%\begin{lstlisting} 
%`\href{https://github.com/inv-ack/inv-ack/blob/7270e64a2600b771f2b1b1b151f7d13fb2ae6c97/repeater.v#L67-L71}{Fixpoint hyperop}` (a n b : nat) : nat :=
%  match n with
%  | 0    => 1 + b
%  | S n' => repeater_from (hyperop a n') (hyperop_init a n') b
%  end.
%\end{lstlisting}
% Edited by Linh
\begin{lstlisting} 
`\href{https://github.com/inv-ack/inv-ack/blob/7270e64a2600b771f2b1b1b151f7d13fb2ae6c97/repeater.v#L67-L71}{\color{blue}Fixpoint hyperop}` (a n b : nat) : nat :=
  match n with
  | 0    => 1 + b
  | S n' => repeater_from
              (hyperop a n') (hyperop_init a n') b
  end.
\end{lstlisting}

\pagebreak
\begin{equation*}
%\begin{array}{l}
A(n,m) ~ \triangleq ~ \begin{cases}
m + 1 & \text{when } n = 0 \\
\rf{\mathcal{A}_{n-1}}{A(n-1,1)\ }(m) & \text{otherwise}
\end{cases} \qquad \qquad \qquad \qquad \qquad \qquad \qquad ~ 
%\end{array}
\end{equation*}
% Linked by A
%\begin{lstlisting}
%`\href{https://github.com/inv-ack/inv-ack/blob/7270e64a2600b771f2b1b1b151f7d13fb2ae6c97/repeater.v#L135-L139}{Fixpoint ackermann}` (n m : nat) : nat :=
%  match n with
%  | 0    => S m
%  | S n' => repeater_from (ackermann n') (ackermann n' 1) m
%  end.
%\end{lstlisting}
% Edited by Linh
\begin{lstlisting}
`\href{https://github.com/inv-ack/inv-ack/blob/7270e64a2600b771f2b1b1b151f7d13fb2ae6c97/repeater.v#L135-L139}{\color{blue}Fixpoint ackermann}` (n m : nat) : nat :=
  match n with
  | 0    => S m
  | S n' => repeater_from
              (ackermann n') (ackermann n' 1) m
  end.
\end{lstlisting}
In the rest of this paper we construct efficient inverses to these
functions.  Our key idea is an inverse to the higher-order \emph{repeater} function, which we call \emph{countdown}.

% Aquinas' promise:
% We explain our core techniques of repeaters and countdowns that allow us to define each level of the Ackermann hierarchy—and their upper inverses—in a straightforward and uniform manner. We show how countdowns, in particular, can be written structurally recursively.

% Anshuman proposes:
% We introduce our core techniques of repeaters and countdowns.
% We show how countdowns, in particular, can be written structurally recursively.










\section*{Roadmap}
\begin{frame}
\frametitle{Roadmap}
\setlength{\leftmargini}{3.8em}
\setbeamertemplate{enumerate items}[default]
\setbeamercolor{local structure}{fg=black}

\imppar{\textbf{Goal.} Inverting $\Ack$ - without computing $\Ack$.}

\bigskip

\begin{enumerate}[\bfseries Step 1.]
	\itemsep 2ex
	\item<done@1->
	Build the hyperoperations/Ackermann hierarchies via \textbf{Repeater}.
	
	\item<come@2->
	Build inverses of the hyperoperations/Ackermann hierarchies via \textbf{Countdown}, the inverse of \emph{Repeater}.
	
	\item
	Use the hierarchy to implement the inverse Ackermann function \\ for inputs encoded in unary. $\bigO (n)$.
	
	\item [\textbf{Bonus.}]
	Improve efficiency for inputs encoded in binary. $\bigO (\log_2n)$.	
\end{enumerate}
\end{frame}


\section{Step 2: Inverting the Hierarchies via Countdown}
\label{sec: countdown}
Many functions on $\mathbb{R}$ are bijections and thus have an intuituve inverse.  
Functions on~$\mathbb{N}$ are often non-bijections and so their inverses
do not come as naturally.

\subsection{Upper Inverses, Expansions, and Repeatability}
%; indeed, simply defining an inverse can be a little subtle.
\begin{defn} \label{defn: inverse}
The \emph{upper inverse} of $F$, written $F^{-1}_{+}$, 
% {\color{red} $F^{-1}$, $F^{-1}_{\mathit{\shortuparrow}}$, $F^{-1}_{\upharpoonleft}$} 
is $\min\{m : F(m)\ge n\}$.\end{defn}
Notice that this is well-defined as long as $F$ is unbounded, \emph{i.e.} $\forall b.~\exists a.~ b \leq F(a)$.  However, as a notion of ``inverse,'' it really only makes sense if $F$ is strictly
increasing, \emph{i.e.} $\forall n,m.~n < m \Rightarrow F(n) < F(m)$, which is, in some sense, the analogue of injectivity in the discrete domain.

We call this function the ``upper inverse'' because, for strictly increasing functions like
addition, multiplication, and exponentiation, the upper inverse is the ceiling of the 
corresponding inverse functions on $\mathbb{R}$. We can characterize inverses more meaningfully as follows: 
\hide{
\color{red} It is reasonable to wonder about the floor.
\begin{defn} \label{defn: lower_inverse}
The converse \emph{lower inverse}, written $F^{-1}_{-}$,
is defined as $\max\{m : F(m)\le n\}$.
\end{defn}
Even if $F$ strictly increases, as $a[n]$ does for every $a\ge 2$, notice that the lower inverse will be undefined for $n < F(0)$, \emph{e.g.} $\{m : a[n]m \le 0 \} = \varnothing$ for $n\ge 3$.
Thus we focus on upper inverses (hereafter just ``inverses''), and discuss lower inverses in \cref{sec: discussion}. 
}%end hide 
\begin{thm} \label{thm: upp-inverse-rel}
	If $F:\mathbb{N}\to \mathbb{N}$ is increasing, then $f$ is the upper inverse of $F$ if and only if $\ \forall n, m.~ f(n)\le m \iff n \le F(m)$.
\end{thm}
\begin{proof}
Fix $n$. The sentence $\forall m.~ f(n)\le m \iff n\le F(m)$ implies: (1) $f(n)$ is a lower bound to $\{m: n \le F(m)\}$ and (2) $f(n)$ is itself in the set, since plugging in $m \triangleq f(n)$ yields $n\le F(f(n))$, which makes $f$ the upper inverse of $F$. Conversely, if $f$ is the upper inverse of $F$, we know $\forall m.~n\le F(m)\Rightarrow f(n)\le m$. Now, $\forall m \ge f(n)$.~$F(m)\ge F(f(n)) \ge n$ by increasing-ness, thus completing the proof.
\end{proof}
\begin{col}
If $F$ is strictly increasing, then $F^{-1}_{+} \circ F$ is
the identity function.
\end{col}
\begin{proof}
By ($\Leftarrow$) of Theorem~\ref{thm: upp-inverse-rel}, $F(n) \le F(n)$ implies 
$(F^{-1}_{+} \circ F)(n) \le n$.  By ($\Rightarrow$) of the same theorem, $(F^{-1}_{+} \circ F)(n) \le (F^{-1}_{+} \circ F)(n)$ implies $F(n) \le F \big((F^{-1}_{+} \circ F)(n)\big)$. $F$ is strictly increasing, so $n \le (F^{-1}_{+} \circ F)(n)$.
\end{proof}

Our setup for inverse requires increasing functions, and our definitions for 
hyperoperations/Ackermann use repeater.  Suppose $F$ is a strictly increasing function.
For a given $a$, is $\rf{F}{a}$ strictly increasing?  No!  For example, the identity function
$\li{id}$ is strictly increasing, but $\rf{\li{id}}{a}(n) = (\li{id} \circ \ldots \circ \li{id}) (a) = a$ is a constant function.  We need a little more.
\begin{defn}
A function $F:\mathbb{N}\to\mathbb{N}$ is an \emph{expansion} if $\forall n.~ F(n)\ge n$. Further, given $a\in \mathbb{N}$, an expansion $F$ is \emph{strict from} $a$ if \note{$\forall n \ge a.~ F(n)\ge n+1$}.
\end{defn}
If $a\ge 1$ and $F$ is an expansion \emph{strict from} $a$, we quickly get: 
$\forall n.~ \rf{F}{a}(n) = F^{(n)}(a) \ge a + n \ge 1 + n$. That is, $\rf{F}{a} \ $ is itself an expansion strict from $0$. 

\begin{defn} \label{rem: repeatable-subset}
$\forall a\ge 1$, if~$F$ is a strictly increasing function 
that is also a strict expansion from~$a$, then $\rf{F}{a}$ preserves 
repeatability.
We thus say that such a function $F$ is \emph{repeatable from} $a$, and denote the 
set of functions repeatable from $a$ as $\repeatable_a$. 
It is straightforward to see that $\forall s,t.~ s \le t \Rightarrow \repeatable_s \subseteq \repeatable_t $.
\end{defn}

\subsection{Contractions and the countdown operation}

Suppose $F \in \repeatable_a$ for some $a \ge 1$, and 
let $f \triangleq F^{-1}_{+}$, \emph{i.e.} the inverse of $F$.
Our goal is to use $f$ to compute an inverse to $\rf{F}{a}$.  
From the preceeding discussion we know that this inverse must exist, 
since $F \in \repeatable_a$ implies $\rf{F}{a} \in \repeatable_0$.  
For reasons that will be clear momentarily, we write this inverse as $\cdt{f}{a}$.  Now
fix $n$ and observe that $\forall m.~f^{(m)}(n)\le a \iff m \ge \cdt{f}{a}(n)$ since
\begin{equation} \label{eq: rf-upp-inv}
\begin{aligned}
\cdt{f}{a}(n)\le m & \iff n\le \rf{F}{a}(m) = F^{(m)}(a) \iff f(n)\le F^{(m-1)}(a) \\
& \iff f^{(2)}(n)\le F^{(m-2)}(a) \iff \ldots \iff f^{(m)}(n)\le a
\end{aligned}
\end{equation}
Moreover, setting $m = \cdt{f}{a}(n)$, we realize that 
$f^{(\cdt{f}{a}(n))}(n) \le a$.  
\textbf{Together these imply that $\cdt{f}{a}(n)$ is the minimum number of 
times~$f$ needs to be compositionally applied to $n$ before equalling or 
passing $a$.} 
In other words, count the length of the chain $\{n, f(n), f^{(2)}(n), \ldots\}$ that 
terminates as soon as we reach/pass $a$.  For this process to work, 
we need each chain link
to be strictly less than the previous, \emph{i.e.} $f$ must be a \emph{contraction}.
\begin{defn} \label{defn: contracting}
	A function $f : \mathbb{N} \to \mathbb{N}$ is a \emph{contraction} if $\forall n.~ f(n) \le n$. Given an $a \ge 1$, a contraction $f$ is 
	\emph{strict above} $a$ if $\forall n > a.~n > f(n)$. We denote the set of contractions by $\contract$ and the set of contractions strict above $a$ by $\contract_a$. Analogously to our observation in 
	Definition \ref{rem: repeatable-subset}, $\forall s\le t.~ \contract_s \subseteq \contract_t$.
\end{defn}
What kinds of functions have contractive inverses? Expansions, naturally:
\begin{thm} \label{thm: expansion-inv-contraction}
$\forall a\in \mathbb{N}.~F\in \repeatable_a \Rightarrow F^{-1}_+ \in \contract_a$.
\end{thm}
\begin{proof}
$\forall n.~F(n)\ge n \Rightarrow n \ge F^{-1}_+(n)$, so $F^{-1}_+ \in
\contract$. Note, $n > a \iff n-1\ge a$.
Since $F\in \repeatable_a$, $F(n-1)\ge n$ holds. 
Next, $n-1\ge F^{-1}_+(n) \Rightarrow n > F^{-1}_+(n)$. 
\end{proof}
This clarifies the inverse relationship between expansions strict from some $a$ and contractions strict above that same $a$. The inverse of an expansion's 
repeater exists and can be built \note{from its own inverse}, in a method formalized as \emph{countdown}.
\begin{defn} \label{defn: informal-countdown} \label{eq: countdown}
Let $f\in \contract_a$. The \textit{countdown to} $a$ of $f$, written 
$\cdt{f}{a}(n)$, is the smallest number of times $f$ needs to be applied to 
$n$ for the answer to equal or go below $a$. \emph{i.e.}, 
$\cdt{f}{a}(n) \triangleq \min\{m: f^{(m)}(n)\le a \}$.
\end{defn}
Inspired by Equation~\ref{eq: rf-upp-inv}, we provide a neat, algebraically manipulable logical sentence equivalent to Definition~\ref{eq: countdown}, which is more useful later in our paper:
\begin{col} \label{col: cdt-repeat}
If $a \in \mathbb{N}$ and $f\in \contract_{a}$, then $\forall n, m.~ \cdt{f}{a}(n)\le m \iff f^{(m)}(n)\le a$.
\end{col}
\begin{proof}
	Fix $a$ and $n$. The interesting direction is $(\Rightarrow)$. Suppose $\cdt{f}{a}(n)\le m$, we get $f^{(m)}(n)\le f^{(\cdt{f}{a}(n))}(n)$ due to $f\in \contract$, and $f^{(\cdt{f}{a}(n))}(n)\le a$ due to Definition~\ref{eq: countdown}.
\end{proof}
%\begin{thm} \label{thm: upp-inv-cdt-rf}
%	For all $a\in \mathbb{N}$, if $f\in \contract_{a}$ is the upper inverse of $F: \mathbb{N}\to \mathbb{N}$, then $\cdt{f}{a}$ is the upper inverse of $\rf{F}{a}$.
%\end{thm}
Another useful result is the recursive formula for \emph{countdown}:
\begin{thm} \label{thm: cdt-recursion}
	For all $a\in \mathbb{N}$ and $f\in \contract_{a}$, $\cdt{f}{a}$ satisfies:
	\begin{equation*}
	\cdt{f}{a}(n) = \begin{cases}
	0 & \text{ if } n \le a \\ 1 + \cdt{f}{a}(f(n)) & \text{ if } n \ge a + 1
	\end{cases}
	\end{equation*}
\end{thm}
\begin{proof}
When $n \le a$, use Corollary~\ref{col: cdt-repeat} as follows: 
$n = f^{(0)}(n)\le a \iff \cdt{f}{a}(n)\le 0$. 
When $n\ge a+1$, proceed by antisymmetry. 
Define $m \triangleq \cdt{f}{a}(f(n))$, and note that 
Corollary~\ref{col: cdt-repeat} gives 
$\cdt{f}{a}(n)\le 1+m \iff f^{(1+m)}(n) \le a$.
Next, a simple expansion of the last clause gives $f^{(m)}(f(n)) \le a$,
and unfolding the definition of~$m$ shows that this last clause is true.
Now since $n\ge a+1$, we have $\cdt{f}{a}(n)\ge 1$ by the above. 
Define $p \triangleq \cdt{f}{a}(n) - 1$. 
It remains to show $\cdt{f}{a}(f(n))\le p \Rightarrow f^{(p)}(f(n))\le a \Rightarrow f^{(p+1)}(n)\le a$, which holds by unfolding $p$'s definition.
\end{proof}

%\begin{thm} \label{thm: cdt-contr-0}
%	For all $a\ge 1$, if $f\in \contract_{a}$, then $\cdt{f}{a}\ \in \contract_{c} \ \forall c$.
%\end{thm}
%\begin{proof}
%	Firstly we show that $\cdt{f}{a}$ is a contraction, namely showing $\cdt{f}{a}(n)\le n \ \forall n$, which has already been proved in the proof of \cref{thm: cdt-repeat}. To show $\cdt{f}{a}$ is strict from $1+c$, it suffices to show it is strict from $1$, equivalently $\cdt{f}{a}(n) \le n - 1 \ \forall n\ge 1$. By \cref{thm: cdt-repeat}, we need to show $f^{(n-1)}(n)\le a$. Assume the converse, then:
%	$$ n \ge 1 + f(n)\ge 2 + f(f(n)) \ge \cdots \ge (n-1) + f^{(n-1)}(n) $$
%	Thus $f^{(n-1)}(n)\le 1 \le a$, a contradiction. The theorem then follows.
%\end{proof}

\subsection{A Structurally Recursive Computation for Countdown}

The higher-order repeater function is well-defined for any input functions, 
even those not in $\repeatable_a$ (although for such functions it may not 
be useful), and so is easy to define
in Coq as shown in \S\ref{sec: countdown-repeater}. In contrast, a 
\emph{countdown} only exists for certain functions, most conveniently 
contractions, which makes it a little harder to encode into Coq. 
Our strategy is to define a \emph{countdown worker} which is written with only
structural recursion and is thus more palatable to Coq, and then 
prove that this worker computes the countdown when passed a contraction.

\begin{defn} \label{defn: countdown-worker}
For any $a\in \mathbb{N}$ and $f: \mathbb{N}\to \mathbb{N}$, the \emph{countdown worker} to $a$ of $f$ is a function $\cdw{f}{a}\ : \mathbb{N}^2\to \mathbb{N}$ such that:
\begin{equation*}
\cdw{f}{a}(n, b) = \begin{cases}
0 & \text{if } b = 0 \vee n\le a \\ 1 +\cdw{f}{a}(f(n), b-1) & \text{if } b \ge 1 \wedge n > a
\end{cases}
\end{equation*}
\end{defn}
The \emph{countdown worker} operates on two arguments: 
the \emph{true argument} $n$, for which we want to simulate 
\emph{countdown to} $a$, 
and the \emph{budget} $b$, the maximum number of times we shall attempt 
to compositionally apply $f$ on the input before giving up. 
If the input goes below or equal $a$ after $k$ applications, \emph{i.e.} $f^{(k)}(n) \le a$, we return the count $k$. If the budget is exhausted (\emph{i.e.} $b = 0$) while the result is still above $a$, we fail by returning the original budget. This definition is admittedly inelegant, but it can clearly be written as a Coq \li{Fixpoint}:
\begin{lstlisting}
`\href{https://github.com/inv-ack/inv-ack/blob/6099297c6ab0e16d14b037fb5ed600c4d22818f6/countdown.v#L94-L100}{Fixpoint countdown\_worker}` (a : nat) (f : nat -> nat) (n k : nat) : nat :=
match k with
 | 0    => 0
 | S k' => match (n - a) with
           | 0 => 0
           | _ => S (countdown_worker a f (f n) k') 
           end
 end.
\end{lstlisting}
Given an $f$ that is a contraction strict from $a$, 
and given a sufficient budget, \emph{countdown worker} 
will compute the correct \emph{countdown} value.  
Careful readers may have noticed that \emph{budget} is similar to 
\emph{gas}, which we discussed in ~\S\ref{sec:incfuncinv} 
and dismissed for potentially 
being too computationally expensive to calculate. 
Budget is actually a refinement of gas because 
we support its use with a method to calculate it in constant time.
In fact, we will soon show that a budget of $n$ is sufficient. 
This lets us define \emph{countdown} in Coq for the first time:
\begin{defn} \label{defn: countdown}
Redefine $\cdt{f}{a}(n) \triangleq \cdw{f}{a}(n, n)$.
\begin{lstlisting}
`\href{https://github.com/inv-ack/inv-ack/blob/6099297c6ab0e16d14b037fb5ed600c4d22818f6/countdown.v#L104}{Definition countdown\_to}` a f n := countdown_worker a f n n.
\end{lstlisting}
\end{defn}
%Before beginning, let us clarify that the definition of $\mathbb{N}$ and operations on $\mathbb{N}$ in Gallina follow the Presburger Arithmetic \cite{presburger}, which despite being weaker than Peano Arithmetic, is a decidable theory. The Coq standard library includes \texttt{Omega}, an extensive listing of provable facts about $\mathbb{N}$ in Presburger Arithmetic, including everything used in this paper, most notably the law of excluded middle for comparisons on $\mathbb{N}$:
%\begin{equation*}
%(n \le m) \vee (m + 1 \le n) \ \ \forall n, m
%\end{equation*}
%, which is provable without the actual law of excluded middle in classical logic. This enables us to prove all results in this paper with Coq's baseline intuitionistic logic. Readers can refer to the \Cref{appendix} for Coq versions of the proofs.
%in which the most operations agree with usual operations on $\mathbb{Z}$, except subtraction, which is defined as:
%\begin{lstlisting}
%Fixpoint sub (n m : nat) : nat :=
%match n with
%| 0 => n
%| S k => match m with
%| 0 => n
%| S l => sub k l end
%end.
%\end{lstlisting}
%Essentially $\li{sub} \ n \ m = \max\{n - m, 0\}$. We will use this subtraction for the rest of the paper.

%Firstly, we begin with a lemma asserting the existence of the countdown value itself. Although its existence is guaranteed by the well-ordering principle of $\mathbb{N}$, we will achieve better by proving it in intuitionistic logic.
%
%\begin{lem} \label{lem: contract-repeat-threshold}
%	For all $a, n\in\mathbb{N}$ and $f\in \contract_{a}$,
%	\begin{equation}
%	\exists m : \left(f^{(m)}(n) \le a \right) \wedge \left(f^{(l)}(n)\le a \Rightarrow m \le l \ \forall l \right)
%	\end{equation}
%\end{lem}
%\begin{proof}
%  Fix $n$ and observe that if $n\le a$, $m = 0$ is the desired choice since $ f^{(0)}(n) = n \le a \ \text{ and } \ 0 \le l \ \forall l $.
%	Consider only when $a\le n$, we can define $c$ such that $n = a + c$. We prove the following statement by induction:
%	\begin{equation*}
%	P(c) \triangleq \exists m : \left(f^{(m)}(n) \le n - c \right) \wedge \left(f^{(l)}(n)\le n - c \Rightarrow m \le l \ \forall l \right)
%	\end{equation*}
%	under assumptions $f\in \contract_{n-c}$ and $c\le n$.
%	\begin{enumerate}[leftmargin=*]
%		\item \textit{Base case.} The case $c = 0$ implies $n = a$, which has been proven above.
%		\item \textit{Inductive step.} Suppose $P(c)$ is proved with witness $m_c$. Note that the assumptions are now $f\in \contract_{n-c}$ and $c+1\le n$, there are two cases:
%		\begin{itemize}[leftmargin=*, label={--}]
%			\item $f^{\left(m_c\right)}(n) = n - c$. Then $f^{\left(m_c+1\right)}(n) \le n - c - 1$. Let $m_{c+1} = m_c + 1$, for all $l$:
%			\begin{equation*}
%			f^{(l)}(n)\le n - c - 1 < f^{\left(m_c\right)}(n) \Rightarrow l > m_c \Rightarrow l \ge m_{c+1}
%			\end{equation*}
%			\item $f^{\left(m_c\right)}(n) \le n - c - 1$. Let $m_{c+1} = m_c$, for all $l$:
%			\begin{equation*}
%			f^{(l)}(n)\le n - c - 1\le n - c \overset{P(c)}{\Rightarrow} l\ge m_{c} = m_{c+1}
%			\end{equation*}
%		\end{itemize}
%		In any cases, we can find a witness $m_{c+1}$ for $P(c+1)$. Thus the proof is complete by induction.\vspace*{-\baselineskip}
%	\end{enumerate}
%\end{proof}

\begin{lem} \label{lem: cdt-init}
	$\forall n\le a.~\cdw{f}{a} (n, b) = 0$.
\end{lem}
\begin{proof}
This simple fact follows directly from Definition \ref{defn: countdown-worker}.
\end{proof}
Next, we show the internal working of \emph{countdown worker} at the $\text{i}^\text{th}$ recursive step, including the accumulated result $1+i$, the current input $f^{(1+i)}(n)$, and the current budget $b-i-1$.
\begin{lem} \label{lem: cdt-intermediate}
	$\forall a, n, b, i \in \mathbb{N}.~\forall f \in \contract$. such that $i < b$ and $a < f^{(i)}(n)$:~
	\begin{equation}  \label{eq: cdt-intermediate}
	\cdw{f}{a}(n, b) = 1 + i + \cdw{f}{a}\left(f^{1+i}(n), b - i - 1\right)
	\end{equation}
\end{lem}
\begin{proof}
	Fix $a$. We proceed by induction on $i$. Then define $P(i)$ as
	\begin{equation*}
	P(i) \triangleq \forall b, n.~ b\ge i+1 \Rightarrow f^{(i)}(n) > a \Rightarrow \cdw{f}{a}(n, b) = 1 + i + \cdw{f}{a}\left(f^{1+i}(n), b - i - 1\right)
	\end{equation*}
	\begin{enumerate}[leftmargin=*]
		\item \textit{Base case.} For $i = 0$, our goal $P(0)$ is:
		$\cdw{f}{a}(n, b) = 1 + \cdw{f}{a}\left(f(n), b - 1\right)$
		where $b \ge 1$ and $f(n)\ge a+1$. This is straightforward.
		\item \textit{Inductive step.} Assume $P(i)$ has been proved. Then $P(i+1)$ is
		\begin{equation*}
		\cdw{f}{a}(n, b) = 2 + i + \cdw{f}{a}\left( f^{2+i}(n), b - i - 2 \right)
		\end{equation*}
		where $b \ge i+2$ and $f^{1+i}(n) \ge a+1$. This also implies $b\ge i+1$ and $\displaystyle f^{(i)}(n) \ge f^{1+i}(n)\ge a+1$ by $f\in \contract$, and so $P(i)$ holds. It suffices to prove:
		\begin{equation*}
		\cdw{f}{a}\left(f^{1+i}(n), b-i-1\right) = 1 + \cdw{f}{a}\left( f^{2+i}(n), b-i-2 \right)
		\end{equation*}
		This is in fact $P(0)$ with $(b, n)$ substituted for $\left(b-i-1, f^{(1+i)}(n)\right)$. Since $f^{(1+i)}(n) \ge a+1$ and $b-i-1\ge 1$, the above holds and $P(i+1)$ follows.
	\end{enumerate}
\end{proof}
Now it is time to prove the correctness of \emph{countdown}, \emph{i.e.} that our computational Definition~\ref{defn: countdown} computes the value originally specified in Definition~\ref{defn: informal-countdown}.
\begin{thm} \label{thm: cdt-repeat}
	$\forall a \in \mathbb{N}.~ \forall f\in \contract_{a}$, we have 
%	\begin{equation} \label{eq: cdt-minimum}
$	\forall n.~ \cdt{f}{a}(n) = \min\left\{ i : f^{(i)}(n) \le a \right\} $.
%	\end{equation}
%	Or equivalently,
%	\begin{equation} \label{eq: cdt-repeat}
%	\cdt{f}{a}(n) \le k \iff f^{(k)}(n) \le a \ \ \forall n, k
%	\end{equation}
\end{thm}
\begin{proof}
%	First, to see why \eqref{eq: cdt-minimum} and \eqref{eq: cdt-repeat} are equivalent, we rewrite \eqref{eq: cdt-minimum} in the following way:
%	$$ \left(f^{(\cdt{f}{a}(n))}(n) \le a\right) \wedge \left(f^{(l)}(n) \le a \Rightarrow \cdt{f}{a}(n) \le l \ \ \forall l\right) \ \ \forall n$$
%	To prove $\eqref{eq: cdt-minimum} \Rightarrow \eqref{eq: cdt-repeat}$, it suffices to show
%	$$ \cdt{f}{a}(n) \le l \Rightarrow f^{(l)}(n) \le a \ \ \forall l $$
%	, which holds due to the fact $\displaystyle f^{(l)}(n) \le f^{(\cdt{f}{a}(n))}(n) \le a$ by $f\in \contract$. To prove $\eqref{eq: cdt-repeat}\Rightarrow \eqref{eq: cdt-minimum}$, it suffices to show $\displaystyle f^{(\cdt{f}{a}(n))}(n) \le a$
%	, which in turn holds by substituting $k$ by $\cdt{f}{a}(n)$ in \eqref{eq: cdt-repeat}. Thus $\eqref{eq: cdt-minimum}\iff \eqref{eq: cdt-repeat}$ and we need only to prove \eqref{eq: cdt-minimum}.
Since $f\in \contract_{a}$ and $\mathbb{N}$ is well-ordered, 
let $m = \min\big\{i : f^{(i)}(n)\le a\big\}$.\footnote{We prove the existence 
of the min in Coq’s intuitionistic logic \href{https://github.com/inv-ack/inv-ack/blob/6099297c6ab0e16d14b037fb5ed600c4d22818f6/countdown.v\#L125-L150}{here} in our codebase.} Our goal becomes 
$\cdt{f}{a}(n) = m$. Note that this setup gives us: 
	\begin{equation}
	\left(f^{(m)}(n) \le a\right) \label{eq: cdt-repeat-tmp} \wedge
	 \left(\forall k.~f^{(k)}(n)\le a \Rightarrow m \le k\right)
	\end{equation}
If $m = 0$, then $n = f^{(0)}(n)\le a$. So $\cdt{f}{a}(n) = \cdw{f}{a}(n, n) = 0 = m$ by Lemma~\ref{lem: cdt-init}.
When $m > 0$, our plan is to apply Lemma~\ref{lem: cdt-intermediate} to get
	\begin{equation*}
	\cdt{f}{a}(n) = \cdw{f}{a}(n, n) = m + \cdw{f}{a}\left(f^{(m)}(n), n-m\right),
	\end{equation*}
and then use Lemma~\ref{lem: cdt-init} over (\ref{eq: cdt-repeat-tmp})'s first conjunct to conclude $\cdt{f}{a}(n) = m$. It suffices to prove the premises of Lemma~\ref{lem: cdt-intermediate}: $a < f^{(m-1)}(n)$ and $m-1 < n$.
	
	The former follows by contradiction: if $f^{(m-1)}(n) \le a$, (\ref{eq: cdt-repeat-tmp})'s second conjunct implies $m\le m-1$, which is impossible for $m > 0$. The latter then easily follows by $f\in \contract_{a}$, since
	$n \ge 1 + f(n) \ge 2 + f(f(n)) \ge \cdots \ge m + f^{(m)}(n)$.
	\linebreak Therefore, $\cdt{f}{a}(n) = m$ in all cases, which completes the proof.
\end{proof}
Theorem~\ref{thm: cdt-repeat} and (\ref{eq: rf-upp-inv}) establish the correctness of the Coq definitions of \emph{countdown worker} and \emph{countdown}, thereby justifying our budget of $n$ and our unification of 
Definitions \ref{defn: informal-countdown} and \ref{defn: countdown}. We wrap everything together with the following theorem:
\begin{thm} \label{thm: cdt-inv-rf}
	$\forall F\in \repeatable_a.~f\triangleq F^{-1}_+$ satisfies $f \in \contract_{a}$ and $\displaystyle \cdt{f}{a} \ = \left(\rf{F}{a}\ \right)^{-1}_+$. Furthermore, if $a\ge 1$, then $\rf{F}{a}\ \in \repeatable_0$ and $\cdt{f}{a}\ \in \contract_0$.
\end{thm}
\begin{proof}
	By Theorem~\ref{thm: expansion-inv-contraction}, $f\triangleq F^{-1}_+ \in \contract_a$. 
	Equation~\ref{eq: rf-upp-inv} and Corollary~\ref{col: cdt-repeat}
	then show $\cdt{f}{a}\ = \left(\rf{F}{a}\ \right)^{-1}_+$.
	Now if $a\ge 1$, a simple induction shows that $F^{(n)}(a)\ge a + n\ge 1 + n$, so $\rf{F}{a}\ \in \repeatable_0$, hence $\cdt{f}{a} \ = \left(\rf{F}{a}\ \right)^{-1}_+ \in \contract_0$ by Theorem~\ref{thm: expansion-inv-contraction}.
\end{proof}



%\section{Inverting Hyperoperations and Ackermann}
%\label{sec: inv-hyperop}
%We now use \emph{countdown} to define the inverse
hyperoperation hierarchy, which features elegant new definitions of
division, $\log$, and $\log^{*}$.
We then modify this technique to arrive at the inverse
Ackermann hierarchy.

\subsection{Inverse hyperoperations, including \li{div}, \li{log}, \li{log*}}

\begin{defn} \label{defn: inv-hyperop}
	The inverse hyperoperations, written $a\angle{n}b$, are defined as:
%	\begin{equation}
%	a\angle{n}b \triangleq \begin{cases}
%	b - 1 & \text{if } n = 0 \\
%	\cdt{a\angle{n-1}}{a_n}(b) & \text{if } n \ge 1
%	\end{cases}
%	\ \ \text{ where } \ a_n = \begin{cases}
%	a & \hspace{-10pt}\text{ if } n = 1 \\
%	0 & \hspace{-10pt}\text{ if } n = 2 \\
%	1 & \hspace{-10pt}\text{ if } n \ge 3
%	\end{cases}
%	\end{equation}
% Edited by Linh
	\begin{equation}
  a\angle{n}b \triangleq \begin{cases}
  b - 1 & \hspace{-5pt}\text{if } n = 0 \\
  \cdt{a\angle{n-1}}{a_n}(b) & \hspace{-5pt}\text{if } n \ge 1
  \end{cases}
   \ \text{ where } \ a_n = \begin{cases}
   a & \hspace{-10pt}\text{ if } n = 1 \\
   0 & \hspace{-10pt}\text{ if } n = 2 \\
  1 & \hspace{-10pt}\text{ if } n \ge 3
  \end{cases}
\end{equation}
% linked by A
%\begin{lstlisting}
%`\href{https://github.com/inv-ack/inv-ack/blob/7270e64a2600b771f2b1b1b151f7d13fb2ae6c97/applications.v#L28-L33}{Fixpoint inv\_hyperop}` (a n b : nat) : nat :=
%  match n with 0 => b - 1 | S n' =>
%    countdown_to (hyperop_init a n') (inv_hyperop a n') b
%  end.
%\end{lstlisting}
% Edited by Linh
\begin{lstlisting}
`\href{https://github.com/inv-ack/inv-ack/blob/7270e64a2600b771f2b1b1b151f7d13fb2ae6c97/applications.v#L28-L33}{\color{blue}Fixpoint inv\_hyperop}` (a n b : nat) : nat :=
  match n with 0 => b - 1 | S n' =>
    countdown_to
      (hyperop_init a n') (inv_hyperop a n') b
  end.
\end{lstlisting}
\end{defn}
\hide{
Continuing with the notion of upper-inverses we have used thus far,
we aim for the \emph{ceiling} of division and $\log$ on real numbers, $\left\lceil b/a \right\rceil$ and $\left\lceil \log_ab \right\rceil$. The iterated logarithm is formally defined as the minimum number of times $\log$ needs to be iteratively
applied to the input for the result to become less than or equal to $1$.
Thus it is natural to define it using \emph{countdown}.

\begin{defn} \label{defn: divc} \label{defn: logc} \label{defn: log*}
\[
\divc(a, b) \! \triangleq \! \cdt{\left(\lambda b. (b{-}a)\right)}{0}(b) \quad
\logc(a, b) \! \triangleq \! \cdt{\left(\lambda b. \divc(a, b)\right)}{1}(b) \quad
\log^*(a, b) \! \triangleq \! \cdt{\left(\lambda b. \logc(a, b)\right)}{1}(b)
%\begin{array}{c}
%\divc(a, b) \triangleq \cdt{\left(\lambda b. (b-a)\right)}{0}(b) \qquad
%\logc(a, b) \triangleq \cdt{\left(\lambda b. \divc(a, b)\right)}{1}(b) \\
%\log^*(a, b) \triangleq \cdt{\left(\lambda b. \logc(a, b)\right)}{1}(b)
%\end{array}
\]
\begin{lstlisting}
Definition divc a b := countdown_to 0 (fun n => n - a) b.
Definition logc a b := countdown_to 1 (divc a) b.
Definition logstar a b := countdown_to 1 (logc a) b.
\end{lstlisting}
\end{defn}
}%end hide
\hide{Proving $\forall a, b \ge 1,~\divc(a, b) = \left\lceil b/a \right\rceil$ serves
to prove that all three definitions are correct, since that statement implies
the correctness of  $\logc$, which in turn implies the correctness of $\log^*$. We will prove this in \ref{blah},
together with the correctness of the entire hyperoperation hierarchy.
}
\noindent where the Curried $a\angle{n{-}1}$ denotes the single-variable function
$\lambda b.a\angle{n{-}1}b$.
We now show that $a\angle{n}$ is the inverse \lb to $a[n]$.
First note that $a\angle{0}\in \contract_0$.~Then:
% Linked by Linh
\begin{lem}
	\href{https://github.com/inv-ack/inv-ack/blob/7270e64a2600b771f2b1b1b151f7d13fb2ae6c97/applications.v#L48-L62}{\color{blue}\coq}
$\forall a.~\forall b.~a\angle{1}b = b - a$.
\end{lem}
\begin{proof}
Theorem~\ref{thm: cdt-recursion} applies because $a\angle{0} \in \contract_0 \subseteq \contract_{a}$, giving the intermediate step shown below.
Thereafter we have $a\angle{1}b = b - a$ by induction on $b$.
%\vspace{-0.8em}
%\begin{equation*}
%a\angle{1}b \quad = \quad \cdt{\left(a\angle{0}\right)}{a}(b) \quad = \quad \begin{cases}
%0 & \text{ if } b\le a \\ 1 + a\angle{1}(b - 1) & \text{ if } b\ge a+1 \end{cases} \\
%\end{equation*}
% Edited by Linh
\begin{equation*}
a\angle{1}b \ = \ \cdt{\left(a\angle{0}\right)}{a}(b) \ = \ \begin{cases}
0 & \text{ if } b\le a \\ 1 + a\angle{1}(b - 1) & \text{ if } b\ge a+1 \end{cases} \\
\end{equation*}
\end{proof}
\begin{col} \label{col: inv-hyperop-1-contr1}
$\forall a \ge 1, a\angle{1} \in \contract_1$.
\end{col}
\hide{\begin{col} \label{col: inv-hyperop-234}
{\color{magenta}$\forall a \in \mathbb{N}$}:
 {\color{blue} For all $a \in \mathbb{N}$},
\begin{enumerate}
	\item If $a\ge 1$, $a\angle{2}b = \divc(a, b) \ \forall b$.
	\item If $a\ge 2$, $a\angle{3}b = \logc(a, b) \ \forall b$ and $a\angle{4}b = \logstarc(a, b) \ \forall b$.
\end{enumerate}
\end{col}}%end hide
%Now we are ready to prove the correctness of the inverse hyperoperations.
\noindent \textbf{N.B.} $a\angle{n}b$ is a total function, but it is never actually used for $a = 0$ when $n \ge 2$ or for $a=1$ when $n \ge 3$. For the values
we do care about, we have our inverse:
% Linked by Linh
\begin{thm} \label{thm: inv-hyperop-correct}
	\href{https://github.com/inv-ack/inv-ack/blob/7270e64a2600b771f2b1b1b151f7d13fb2ae6c97/applications.v#L72-L98}{\color{blue}\coq}
When $n\le 1$, or $n \le 2 \wedge a\ge 1$, or $a\ge 2$, then
$a\angle{n} = \left(a[n]\right)^{-1}$.
\end{thm}
\begin{proof}
$\forall n \ge 2,~$ let $a_0 = a, a_1 = 0, a_n = 1$. Define $P$ and $Q$ as:
\begin{equation*}
P(n) \triangleq  \Big(a[n] \in \repeatable_{a_n}\Big)\ \text{ and } \
Q(n) \triangleq  \Big((a\angle{n} = \left(a[n]\right)^{-1} \Big).
\end{equation*}
We have three goals:
\begin{enumerate}
	\item $\forall a.~ Q(0) \wedge Q(1)$ \\ \vspace{-1.2em}
	\item $\forall a \ge 1.~ Q(2)$ \\ \vspace{-1.2em}
	\item $\forall a \ge 2.~ Q(n)$ \\ \vspace{-1.2em}
\end{enumerate}
	Note that \ \mbox{$\forall n.~ a\angle{n+1}$} $= \cdt{a\angle{n}}{a_n}$ ~and \ $a[n+1] = \rf{a[n]}{a_n}$. By Theorem~\ref{thm: cdt-inv-rf},
% Edited by Linh
%\vspace*{-1em}
%\noindent\begin{minipage}{.45\linewidth}
%\begin{equation}
%P(n) \Rightarrow Q(n) \Rightarrow Q(n+1) \label{eq: tmp-induction-1}
%\end{equation}
%\end{minipage}
%\begin{minipage}{.55\linewidth}
%\begin{equation}
%\qquad a_{n} \ge 1 \Rightarrow P(n) \Rightarrow P(n+1) \label{eq: tmp-induction-2}
%\end{equation}
%\end{minipage}
% Edited by Linh
	\begin{align}
	P(n) & \ \Rightarrow \ Q(n) \ \Rightarrow \ Q(n+1) \label{eq: tmp-induction-1} \\
	\qquad a_{n} \ge 1 & \ \Rightarrow \ P(n) \ \Rightarrow \ P(n+1) \label{eq: tmp-induction-2}
	\end{align}
Goal 1: $P(0) \iff \lambda b.(b+1)\in \repeatable_a$ and
$Q(0)\hspace{-0.3em}\iff\hspace{-0.3em}a\angle{0}\hspace{-0.3em}=\hspace{-0.3em}\left(a[0]\right)^{-1} \lb \iff (b-~1\le c \iff b\le c+1)$.
These are both straightforward, and $Q(1)$ holds by~(\ref{eq: tmp-induction-1}). 
Goal 2: we have $a \ge 1$, so
$P(1)$ holds by $P(0)$ and~(\ref{eq: tmp-induction-2}), and
$Q(2)$ holds by $Q(1)$ and~(\ref{eq: tmp-induction-1}).
Goal 3: we have $a\ge 2$, and using~(\ref{eq: tmp-induction-1}) and $Q(0)$
reduces the goal to $P(n)$. Using~(\ref{eq: tmp-induction-2}) and the fact that $\forall n \neq 1.~a_n\ge 1$, the goal reduces to $P(2)$. This unfolds to:
\begin{equation*}
a[2]\in \repeatable_0 \iff \forall b < c.~ab < ac \quad \wedge \quad \forall b \ge 1.~ab \ge b+1\text{,}
\end{equation*}
which is straightforward for $a\ge 2$. Induction on $n$ gives us the third goal.
\end{proof}
%With this theorem, we have proved the correctness of $\divc$, $\logc$ and $\log^*$, as well as the whole inverse hyperoperation hierarchy.
\begin{rem}
Three early hyperoperations are $a[2]b = ab$, $a[3]b = a^b$ and
$a[4]b = \! ^ba$, so, by Theorem~\ref{thm: inv-hyperop-correct}, we can define their inverses $\left\lceil b/a \right\rceil$, $\left\lceil \log_a b \right\rceil$, and $\log^*_a b$ as	
\href{https://github.com/inv-ack/inv-ack/blob/7270e64a2600b771f2b1b1b151f7d13fb2ae6c97/applications.v#L102-L113}{\color{blue}$a\angle{2}b$},	
\href{https://github.com/inv-ack/inv-ack/blob/7270e64a2600b771f2b1b1b151f7d13fb2ae6c97/applications.v#L115-L124}{\color{blue}$a\angle{3}b$}, and	
\href{https://github.com/inv-ack/inv-ack/blob/7270e64a2600b771f2b1b1b151f7d13fb2ae6c97/applications.v#L126-L128}{\color{blue}$a\angle{4}b$}.
Note that the functions $\log_a b$ and $\log^*_a b$ are not in the Coq Standard Library but are one-liners for us:

\begin{lstlisting}
Definition divc a b := inv_hyperop a 2 b.
Definition logc a b := inv_hyperop a 3 b.
Definition logstar a b := inv_hyperop a 4 b.
\end{lstlisting}
\end{rem}



\hide{
\begin{col}
	$\forall a \in \mathbb{N}$:
	\begin{enumerate}
		\item If $a\ge 1$, $\divc(a, b) = \left\lceil \frac{b}{a} \right\rceil$.
		\item If $a\ge 2$, $\logc(a, b) = \left\lceil \log_ab \right\rceil $ and $\logstarc(a, b) = \log^*_ab$.
	\end{enumerate}
\end{col}
\begin{proof}
	The result follows from \cref{col: inv-hyperop-234}, \cref{thm: inv-hyperop-correct} and the first few expansions of the hyperoperation hierarchy: $a[2]b = ab$, $a[3]b = a^b$ and $a[4]b = \ ^ba$.
\end{proof}

\begin{rem}
	\Cref{thm: cdt-recursion} leads to useful recursive formulas for $\divc$, $\logc$ and $\logstarc$:
	\begin{align*}
	\divc(a, b) & = 1 + \divc(a, b - a) & \text{ if } b > 0 &\text{, and } 0 \text{ otherwise}. \\
	\logc(a, b) & = 1 + \logc\left(a, \left\lceil b/a \right\rceil\right) & \text{ if } b > 1 &\text{, and } 0 \text{ otherwise}. \\
	\logstarc(a, b) & = 1 + \divc(a, \left\lceil \log_ab \right\rceil) & \text{ if } b > 1 &\text{, and } 0 \text{ otherwise}.
	\end{align*}
\end{rem}
}

\subsection{The inverse Ackermann hierarchy}
\label{subsec: inv_ack_hier}

% NEW METHOD FOR INVERSE ACKERMANN
Next, we want to use \emph{countdown} to build the \emph{inverse Ackermann hierarchy}, where each
level $\alpha_i$ inverts the level $\Ack_i$.
We know $\Ack_{i+1} = \rf{\Ack_i}{\Ack_i(1)}$\hspace{0.2em}, so the recursive
rule $\alpha_{i+1} \triangleq \cdt{\alpha_i}{\Ack_i(1)}$ is tempting.
But this approach is flawed because it still depends on $\Ack_i$.
Instead, we reexamine the inverse relationship: suppose $\alpha_i = \big(\Ack_i\big)^{-1}$ and $\alpha_{i+1} = \big(\Ack_{i+1}\big)^{-1}$. Then $\Ack_{i+1}(m) = \big(\Ack_{i}\big)^{(m)} \big(\Ack_{i}(1)\big)$. We then have:
\begin{equation} \label{eq: inv-ack-hier-derive}
\alpha_{i+1}(n)\le m \iff n\le \big(\mathcal{A}_i\big)^{(m+1)}(1) \iff \big(\alpha_i\big)^{(m+1)}(n) \le 1
\end{equation}
Equivalently, $\alpha_{i+1}(n) = \min\big\{m : \big( \alpha_i \big)^{(m+1)}(n)\le 1\big\}$, or $\alpha_{i+1}(n) = \cdt{\alpha_i}{1}\big(\alpha_i(n)\big)$. From \eqref{eq: inv-ack-hier-derive} we can thus define the inverse Ackermann hierarchy:
\begin{defn} \label{defn: inv-ack-hier}
	$ \alpha_i \triangleq \begin{cases}
	\lambda n.(n - 1) & \text{if } i = 0
	\\ \big(\cdt{\alpha_{i-1}}{1}\big)\circ \alpha_{i-1} & \text{if } i\ge 1 \end{cases}
$
\end{defn}
Abracadabra! We can express $\alpha$ using the hierarchy \emph{without referring to $\Ack$ itself}:
\begin{thm} \label{thm: inv-ack-new}
	%For all $n, k$, $n\le \Ack(k, k) \!\!\iff \!\! \alpha_k(n)\le k$. Thus
	$\forall n.~ \alpha(n) = \min\big\{k : \alpha_k(n)\le k \big\}$.
\end{thm}


\begin{table}[t]
	\begin{centermath}
		\begin{array}{c | c@{\hspace{1.5em}} c@{\hspace{1.5em}} c@{\hspace{1.5em}} c@{\hspace{1.5em}} c@{\hspace{1.5em}} c@{\hspace{1.5em}} c@{\hspace{1.5em}} c@{\hspace{1.5em}} c@{\hspace{1.5em}} c@{\hspace{1.5em}}}
			          $\diagbox[height=\line]{$\alpha_k$}{$n$}$ & 0\tikzmark{zero} & 1\tikzmark{one} & \tikzmark{twoa}2\tikzmark{two} & \tikzmark{threea}3\tikzmark{three} & \tikzmark{foura}4\tikzmark{four} & \tikzmark{fivea}5\tikzmark{five} & \tikzmark{sixa}6\tikzmark{six} & \tikzmark{sevena}7\tikzmark{seven} & \tikzmark{eighta}8\tikzmark{eight} & \tikzmark{ninea}9\tikzmark{nine} \\
			\hline
			\alpha_0 & 0\tikzmark{zero_succ} & 0\tikzmark{one_succ} & \tikzmark{two_fail}1 & \tikzmark{three_fail}2 & \tikzmark{four_fail1}3 & \tikzmark{five_fail1}4 & \tikzmark{six_fail1}5 & \tikzmark{seven_fail1}6 & \tikzmark{eight_fail1}7 & \tikzmark{nine_fail1}8 \\
			\alpha_1 & 0 & 0 & 0\tikzmark{two_succ} & 1\tikzmark{three_succ} & \tikzmark{four_fail2}2 & \tikzmark{five_fail2}3 & \tikzmark{six_fail2}4 & \tikzmark{seven_fail2}5 & \tikzmark{eight_fail2}6 & \tikzmark{nine_fail2}7 \\
			\alpha_2 & 0 & 0 & 0 & 0 & 1\tikzmark{four_succ} & 1\tikzmark{five_succ} & 2\tikzmark{six_succ} & 2\tikzmark{seven_succ} & \tikzmark{eight_fail3}3 & \tikzmark{nine_fail3}3 \\
			\alpha_3 & 0 & 0 & 0 & 0 & 0 & 0 & 1 & 1 & 1\tikzmark{eight_succ} & 1\tikzmark{nine_succ} \\
		\end{array}
	\end{centermath}
	\caption{Intuition for $\alpha(n)$ defined without $\Ack(n)$.}
	\label{table:inv_intuition}
\end{table}
  \begin{tikzpicture}[overlay, remember picture, yshift=.25\baselineskip, shorten >=.5pt, shorten <=.5pt]
    \draw [->, green] ({pic cs:zero}) to [bend left]({pic cs:zero_succ});
    \draw [->, green] ({pic cs:one}) to [bend left]({pic cs:one_succ});
   	\draw [->, red] ({pic cs:twoa}) to [bend right]({pic cs:two_fail});
   	\draw [->, green] ({pic cs:two}) to [bend left]({pic cs:two_succ});
   	\draw [->, red] ({pic cs:threea}) to [bend right]({pic cs:three_fail});
   	\draw [->, green] ({pic cs:three}) to [bend left]({pic cs:three_succ});
   	\draw [->, red] ({pic cs:foura}) to [bend right]({pic cs:four_fail1});
   	\draw [->, red] ({pic cs:foura}) to [bend right]({pic cs:four_fail2});
   	\draw [->, green] ({pic cs:four}) to [bend left]({pic cs:four_succ});
   	\draw [->, red] ({pic cs:fivea}) to [bend right]({pic cs:five_fail1});
   	\draw [->, red] ({pic cs:fivea}) to [bend right]({pic cs:five_fail2});
   	\draw [->, green] ({pic cs:five}) to [bend left]({pic cs:five_succ});
   	\draw [->, red] ({pic cs:sixa}) to [bend right]({pic cs:six_fail1});
   	\draw [->, red] ({pic cs:sixa}) to [bend right]({pic cs:six_fail2});
   	\draw [->, green] ({pic cs:six}) to [bend left]({pic cs:six_succ});
   	\draw [->, red] ({pic cs:sevena}) to [bend right]({pic cs:seven_fail1});
   	\draw [->, red] ({pic cs:sevena}) to [bend right]({pic cs:seven_fail2});
   	\draw [->, green] ({pic cs:seven}) to [bend left]({pic cs:seven_succ});
   	\draw [->, red] ({pic cs:eighta}) to [bend right]({pic cs:eight_fail1});
   	\draw [->, red] ({pic cs:eighta}) to [bend right]({pic cs:eight_fail2});
   	\draw [->, red] ({pic cs:eighta}) to [bend right]({pic cs:eight_fail3});
   	\draw [->, green] ({pic cs:eight}) to [bend left]({pic cs:eight_succ});
   	\draw [->, red] ({pic cs:ninea}) to [bend right]({pic cs:nine_fail1});
   	\draw [->, red] ({pic cs:ninea}) to [bend right]({pic cs:nine_fail2});
   	\draw [->, red] ({pic cs:ninea}) to [bend right]({pic cs:nine_fail3});
   	\draw [->, green] ({pic cs:nine}) to [bend left]({pic cs:nine_succ});
  \end{tikzpicture}

% \vspace{-3em}
Table~\ref{table:inv_intuition} gives intuition for our new definition of $\alpha$.
This table is populated exactly as suggested by 
Table~\ref{table: hyperop-ack-inv}:
$\alpha_{0}(n)$ is $n-1$, 
$\alpha_{1}(n)$ is $n-2$,
$\alpha_{2}(n)$ is $\left\lceil \frac{n-3}{2} \right\rceil$, and
$\alpha_{3}(n)$ is $\left\lceil \log_2 ~ (b + 3)\right\rceil - 3$.
Now for any $n$, we search down its column, which is indexed by $k$, 
looking for the smallest $k$ such that $\alpha_k(n) \le k$.

Consider $n = 8$. $\alpha_{0}(8) = 7$. Because $7 \not\le 0$, $\alpha_0$ is rejected.
Similarly $\alpha_{1}(8) = 6 \not\le 1$ and $\alpha_{2}(8) = 3 \not\le 2$
are rejected. Finally $\alpha_{3}(8) = 1$ is accepted because $1 \le 3$. 
Indeed, $\alpha(8) = 3$.
Similar reasoning applies to the other values for $n$. To aid intuition in this table,
we show show unsuccessful searches in {\color{red}red} and successful searches in 
{\color{green}green}.

Now all that remains is to provide a structurally-recursive function that computes $\alpha$.
%\begin{defn} \label{defn: inv-ack-worker}
%	The inverse Ackermann worker is a function $\alpha^{\W}$: %(\mathbb{N}\to \mathbb{N}) \times \mathbb{N}^3\to \mathbb{N}$ such that for all $n, k, b\in \mathbb{N}$ and $f:\mathbb{N}\to \mathbb{N}$:
%	\begin{equation} \label{eq: inv-ack-worker-recursion}
%	\alpha^{\W}(f, n, k, b) \triangleq \begin{cases}
%	k & \text{if } b = 0 \vee n\le k \\ \alpha^{\W}(\cdt{f}{1}\circ f , \cdt{f}{1}(n), k+1, b-1) & \text{if } b \ge 1 \wedge n \ge k+1
%	\end{cases}
%	\end{equation}
%\end{defn}
% Edited by Linh
\begin{defn} \label{defn: inv-ack-worker}
	The inverse Ackermann worker, written $\alpha^{\W}$, is a function from $\mathbb{N}^4$ to $\mathbb{N}$ defined as: %(\mathbb{N}\to \mathbb{N}) \times \mathbb{N}^3\to \mathbb{N}$ such that for all $n, k, b\in \mathbb{N}$ and $f:\mathbb{N}\to \mathbb{N}$:
	\begin{equation} \label{eq: inv-ack-worker-recursion}
	\begin{aligned}
	& \alpha^{\W}(f, n, k, b) \\ 
	& \triangleq \begin{cases}
	k & \text{if } b = 0 \vee n\le k \\ \alpha^{\W}(\cdt{f}{1}\circ f , \cdt{f}{1}(n), k+1, b-1) & \text{if } b \ge 1 \wedge n \ge k+1
	\end{cases}
		\end{aligned}
	\end{equation}
\end{defn}
\noindent Next, we show that this function computes the inverse Ackermann function when passed appropriate arguments.
\hide{
Given the arguments $\big(\alpha_i, \alpha_i(n), i, b - i\big)$ such that $\alpha_i(n) > i$ and $b > i$, $\alpha^{\W}$ takes on arguments $\big(\alpha_{i+1}, \alpha_{i+1}(n), i+1, b - (i+1)\big)$ at the next recursive call. Thus if $\alpha^{\W}$ is given a sufficient budget $b$, it will recursively transform the tuple $(\alpha_k, \alpha_k(n), k, b - k)$ until a point $k$ where $\alpha_k(n)\le k$, and will then return $k$. We now need to show that $\alpha^{\W}$ correctly computes $\alpha(n)$ given a reasonable budget.
The following theorem demonstrates that a budget of $n$ suffices.
}%end hide
\hide{
\noindent We are finally ready for a strategy to compute the inverse Ackermann function:
}%end hide
%formalizes a setting for $\W\alpha$ to work.
% Linked by Linh
\begin{thm} \label{thm: inv-ack-worker-correct}
	\href{https://github.com/inv-ack/inv-ack/blob/7270e64a2600b771f2b1b1b151f7d13fb2ae6c97/inv_ack.v#L199-L231}{\color{blue}\coq}
	$\forall n.~\alpha^{\W}\big(\alpha_0, \alpha_0(n), 0, n\big) = \alpha(n)$.
\end{thm}
\begin{proof}[Proof outline]
	When given the arguments $\big(\alpha_i, \alpha_i(n), i, b - i\big)$ such that $\alpha_i(n) > i$ and $b > i$, $\alpha^{\W}$ takes on the arguments $\big(\alpha_{i+1}, \alpha_{i+1}(n), i+1, b - (i+1)\big)$ at the next recursive call. A simple induction on $k$ then shows that if $k\le \min\big\{b, \alpha_k(n)\big\}$,
	\begin{equation} \label{eq: inv-ack-worker-intermediate}
	\alpha^{\W}\big(\alpha_0, \alpha_0(n), 0, b\big) = \alpha^{\W}\big(\alpha_k, \alpha_k(n), k, b-k\big)
	\end{equation}
	Let $m \triangleq \min\big\{k : \alpha_k(n) \le k \}$. Then $m\le n$ since $\alpha_n(n)\le n$. \eqref{eq: inv-ack-worker-intermediate} then implies:
	$$ \alpha^{\W}\big(\alpha_0, \alpha_0(n), 0, n\big) = \alpha^{\W}\big(\alpha_m, \alpha_m(n), m, n - m\big) = m = \alpha(n) $$
\end{proof}
\noindent We put a more involved mathematical proof of correctness in Appendix~\ref{apx:proof_correct_inv_ack_worker}, and a mechanized proof
% Linked by Linh
	\href{https://github.com/inv-ack/inv-ack/blob/7270e64a2600b771f2b1b1b151f7d13fb2ae6c97/inv_ack.v#L163-L231}{\color{blue}here}.
We thus have a (re-)definition of inverse Ackermann that is %computationally
definable via a Coq-accepted worker, \emph{i.e.} $\alpha(n) \triangleq \alpha^{\W}\big(\alpha_0, \alpha_0(n), 0, n\big)$. Below, we present the inverse
Ackermann function in Gallina.

% Linked by A
%\begin{lstlisting}
%`\href{https://github.com/inv-ack/inv-ack/blob/7270e64a2600b771f2b1b1b151f7d13fb2ae6c97/inv_ack.v#L155-L161} {Fixpoint inv\_ack\_wkr}` (f : nat -> nat) (n k b : nat) : nat :=
%  match b with 0 => 0 | S b' =>
%    if (n <=? k) then k else let g := (countdown_to f 1) in
%      inv_ack_wkr (compose g f) (g n) (S k) b
%  end.
%
%`\href{https://github.com/inv-ack/inv-ack/blob/7270e64a2600b771f2b1b1b151f7d13fb2ae6c97/inv_ack.v#L37-L41}{Fixpoint alpha}` (m x : nat) : nat :=
%  match m with 0 => x - 1 | S m' =>
%    countdown_to 1 (alpha m') (alpha m' x)
%  end.
%
%`\href{https://github.com/inv-ack/inv-ack/blob/7270e64a2600b771f2b1b1b151f7d13fb2ae6c97/inv_ack.v#L167}{Definition inv\_ack}` := inv_ack_wkr (alpha 0) (alpha 0 n) 0 n.
%\end{lstlisting}
% Edited by Linh
\begin{lstlisting}
`\href{https://github.com/inv-ack/inv-ack/blob/7270e64a2600b771f2b1b1b151f7d13fb2ae6c97/inv_ack.v#L155-L161} {\color{blue}Fixpoint inv\_ack\_wkr}` f n k b :=
  match b with 0 => 0 | S b' =>
    if (n <=? k) then k
      else let g := (countdown_to f 1) in
        inv_ack_wkr (compose g f) (g n) (S k) b
  end.

`\href{https://github.com/inv-ack/inv-ack/blob/7270e64a2600b771f2b1b1b151f7d13fb2ae6c97/inv_ack.v#L37-L41}{\color{blue}Fixpoint alpha}` m x :=
  match m with 0 => x - 1 | S m' =>
    countdown_to 1 (alpha m') (alpha m' x)
  end.

`\href{https://github.com/inv-ack/inv-ack/blob/7270e64a2600b771f2b1b1b151f7d13fb2ae6c97/inv_ack.v#L167}{\color{blue}Definition inv\_ack}` :=
  inv_ack_wkr (alpha 0) (alpha 0 n) 0 n.
\end{lstlisting}

\noindent Note that this is not the linear-time computation we presented in
Figure~\ref{fig:standalone}. We will arrive at that code via an improvement discussed in the next section.
% OLD METHOD FOR INVERSE ACKERMANN

%Using the Ackermann kludge from \cref{sec: overview}, we sketch a method to find the inverse Ackermann function from the inverse hyperoperations.
%\begin{thm} \label{thm: inv-hyperop-inv-ack}
%	For all $n, k$, $n\le \Ack(k, k) \iff 2\angle{k}(n+3)\le k+3$.
%\end{thm}
%\begin{proof}
%	$n\le \Ack(k, k) \iff n\le 2[k](k+3) - 3$. Now $2[k]$ is an expansion by \cref{thm: inv-hyperop-correct}'s proof, so $2[k](k+3) \ge k+3\ge 3$, so  $n\le 2[k](k+3) - 3 \iff n+3 \le 2[k](k+3)\iff 2\angle{k}(n+3)\le k+3$, again by \cref{thm: inv-hyperop-correct}.
%\end{proof}
%\Cref{thm: inv-hyperop-inv-ack} allows a simple method to compute the Inverse Ackermann function, based on the \emph{budget} idea in \emph{countdown worker}.
%\begin{defn} \label{defn: inv-ack-worker}
%	The \emph{inverse Ackermann worker} of $f$ is a function $\W\alpha\ : \mathbb{N}^4\to \mathbb{N}$ such that for all $n, k, b\in \mathbb{N}$ and $f:\mathbb{N}\to \mathbb{N}$:
%	$$ \W\alpha(f, n, k, b) = \begin{cases}
%	k - 3 & \text{if } b = 0 \vee f(n)\le k+3 \\ \W\alpha(\cdt{f}{a_k}\ , n, k+1, b-1) & \text{if } b \ge 1 \wedge f(n) > k+3
%	\end{cases} $$
%	where $a_0 = 2$, $a_1 = 0$ and $a_k = 1 \ \forall k\ge 2$.
%\end{defn}
%The function $\W\alpha$ takes a function $f$ and pretends it was the first level of the inverse hyperoperation hierarchy. It then keeps applying countdown to $f$, supposedly arrives at the $\text{m}^{\text{th}}$-level at the $\text{m}^{\text{th}}$-recursive step, until the budget $b$ is exhausted or $f(n)\le k+3$, i.e. \emph{early stopping}, and will output $k$. If at the beginning $k=3$ and $f = 2\angle{0}$, the early stopping condition becomes $2\angle{k}n \le k+3$, which when replace $n$ by $n+3$ gives what we need in \cref{thm: inv-hyperop-inv-ack}.
%The inverse Ackermann function can be defined as follows.
%\begin{thm} For all $n$, we have $\alpha(n) = \W\alpha(\lambda m.(m - 1), n+3, 3, n)$.
%\end{thm} 

\section*{Roadmap}
\begin{frame}
\frametitle{Roadmap}
\setlength{\leftmargini}{3.8em}
\setbeamertemplate{enumerate items}[default]
\setbeamercolor{local structure}{fg=black}

\imppar{\textbf{Goal.} Inverting $\Ack$ - without computing $\Ack$.}

\bigskip

\begin{enumerate}[\bfseries Step 1.]
	\itemsep 2ex
	\item<done@1->
	Build the hyperoperations/Ackermann hierarchies via \textbf{Repeater}.
	
	\item<done@1->
	Build inverses of the hyperoperations/Ackermann hierarchies via \textbf{Countdown}, the inverse of \emph{Repeater}.
	
	\item<come@2->
	Use the hierarchy to implement the inverse Ackermann function \\ for inputs encoded in unary. $\bigO (n)$.
	
	\item [\textbf{Bonus.}]
	Improve efficiency for inputs encoded in binary. $\bigO (\log_2n)$.	
\end{enumerate}
\end{frame}


\section{Step 3: The Inverse Ackermann function in linear time for unary encoding system}
\label{sec: inv-ack}
In this section, we provide a time analysis for the previously defined computation of the inverse Ackermann function. In particular, we prove that its running time is $O(n^2)$. We then provide a simple improvement that can cut the running time to $O(n\cdot\alpha(n))$, and a subsequent improvement that ultimately reduces it to $O(n)$.
\subsection{Time Analysis}
In this section, let us forget about the axiom of extensionality and identify each function on $\mathbb{N}$ with its \emph{computation}, i.e. the program that computes it in Coq. We will be careful not to conclude $f = g$ when they agree on all inputs but are computed with different pieces of code. It allows us to formalize a definition of running time of functions.
\begin{defn}
	Given a function $f:\mathbb{N}\to\mathbb{N}$ in Coq, the \emph{running time} of $f$ on input $n\in \mathbb{N}$, denoted by $\runtime(f, n)$ is the total number of computational steps it takes to compute $f(n)$.
\end{defn}
Below are several useful lemmas for our analysis.
\begin{lem} \label{lem: sub-runtime}
	For all $a, n$, $\runtime(\lambda n.(n - a), n) = \Theta(\min\{a, n\})$ per subtraction's Coq definition.
\end{lem}
\begin{lem} \label{lem: compose-runtime}
	Per composition's Coq definition, for all $f$ and $g: \mathbb{N}\to \mathbb{N}$, $\runtime(f\circ g, n) = \runtime(f, g(n)) + \runtime(g, n)$.
\end{lem}
\begin{lem} \label{lem: cdt-runtime}
	Per our definition of \emph{countdown} (\Cref{defn: countdown}), for all $a\ge 1$ and $f\in \contract_{a}$, $\runtime\big(\cdt{f}{a}\ , n\big) = 1 \ \forall n\le 1$ and
	\begin{equation*}
	\runtime\big(\cdt{f}{a}\ , n\big) = \sum_{i=0}^{\cdt{f}{a}(n) - 1} \runtime\left(f, f^{(i)}(n)\right) + \Theta\big(a\cdot \cdt{f}{a}(n)\big) \ \ \forall n\ge 2
	\end{equation*}
\end{lem}
From \cref{lem: compose-runtime} and \cref{lem: cdt-runtime}, the following lemma easily follows.
\begin{lem} \label{lem: inv-ack-hier-runtime}
	Per \cref{defn: inv-ack-hier} and \cref{defn: countdown}, for all $i$, 
	\begin{equation*}
	\runtime\big(\alpha_{i+1}, n\big) = \sum_{k=0}^{\alpha_{i+1}(n)}\runtime\left(\alpha_i, \alpha_i^{(k)}(n)\right) + \Theta\big(\alpha_{i+1}(n)\big)
	\end{equation*}
\end{lem}
This lemma implies $\runtime\big(\alpha_{i+1}, n\big) \ge \runtime\big(\alpha_i, n\big)$. If for some $i$, $\runtime\big(\alpha_i, n\big) = \Theta\big(n^2\big)$, each function in the hierarchy will take at least $\Omega\big(n^2\big)$ to compute from $\alpha_i$, thus making $\runtime(\alpha, n) = \Omega\big(n^2\big)$ per \cref{defn: inv-ack-worker}. The next lemma shows $i = 2$ suffices.
\begin{lem}
	Per \cref{defn: inv-ack-hier}, $\runtime\big(\alpha_2, n\big) = \Theta\big(n^2\big) \ \forall n$. 
\end{lem}
\begin{proof}
	$\alpha_1 = \big(\cdt{\lambda m.(m-1)}{1}\big)\circ \big(\lambda m.(m-1)\big) \equiv \lambda m.(m - 2)$, by \cref{lem: inv-ack-hier-runtime},
	\begin{equation*}
	\runtime\big(\alpha_1, n\big) = \sum_{i=0}^{n-1} \runtime\big(\lambda m.(m-1), n - i\big) + \Theta\big((n-2)\big) = \Theta(n)
	\end{equation*}
	, since $\runtime\big(\lambda m.(m-1), k\big) = 1 \ \forall k$. Since $\alpha_2 = \big(\cdt{\alpha_1}{0}\big)\circ \alpha_1 $, again by \cref{lem: inv-ack-hier-runtime},
	\begin{equation*}
	\runtime\big(\alpha_2, n\big)
	= \sum_{i=0}^{\lceil (n-3)/2 \rceil} \runtime \big(\alpha_1, n-2i\big) + \Theta\left(\frac{n}{2}\right)
	= \Theta\left( \sum_{i=0}^{\lceil (n-3)/2 \rceil}(n - 2i) \right)
	= \Theta\big(n^2\big)
	\end{equation*}
	The proof is complete.
\end{proof}
\begin{col}
	$\runtime(\alpha, n) = \Omega\big(n^2\big)$ per \cref{defn: inv-ack-worker}.
\end{col}
Intuitively, the function $\alpha_2 \equiv \lambda n.(n-2)$ is responsible for dragging the whole hierarchy's performance due to one silly weakness: its does not know it will always output $n-2$ before beginning its computation, hence needs to tediously subtract $1$ until it goes below $2$. This observation leads to the next improvement.

%This lemma implies $\runtime\big(\cdt{f}{a}\ , n\big) \ge \runtime(f, n)$. If for some $i$, $\runtime\big(2\angle{i}, n\big) = \Theta\big(n^2\big)$, the entire hierarchy will take at least $\Omega\big(n^2\big)$ from $2\angle{i}$, thus making $\runtime(\alpha, n) = \Omega\big(n^2\big)$ per \cref{defn: inv-ack-worker}. The next lemma shows $i = 2$ suffices.
%\begin{lem}
%	Per \cref{defn: inv-hyperop}, $\runtime\big(2\angle{2}, n\big) = \Theta\big(n^2\big) \ \forall n$. 
%\end{lem}
%\begin{proof}
%	Since $2\angle{1} = \cdt{\lambda m.(m-1)}{2}\ \equiv \lambda m.(m - 2)$, by \cref{lem: cdt-runtime},
%	$$ \runtime\big(2\angle{1}, n\big) = \sum_{i=0}^{n-2} \runtime\big(\lambda m.(m-1), n - i\big) + \Theta\big(3(n-2)\big) = \Theta(n) $$
%	, since $\runtime\big(\lambda m.(m-1), k\big) = 1 \ \forall k$. Since $2\angle{2} = \cdt{2\angle{1}}{0}\ $, again by \cref{lem: cdt-runtime},
%	$$ \runtime\big(2\angle{2}, n\big)
%	= \sum_{i=0}^{\lceil n/2 \rceil} \runtime \big(2\angle{1}, n-2i\big) + \Theta\left(\frac{n}{2}\right)
%	= \Theta\left( \sum_{i=0}^{\lceil n/2 \rceil}(n - 2i) \right)
%	= \Theta\big(n^2\big) $$
%	The proof is complete.
%\end{proof}
%\begin{col}
%	$\runtime(\alpha, n) = \Omega\big(n^2\big)$ per \cref{defn: inv-ack-worker}.
%\end{col}
%Intuitively, the function $2\angle{1} \equiv \lambda n.(n-2)$ is responsible for dragging the whole hierarchy's performance due to one silly weakness: its does not know it will always output $n-2$ before beginning its computation, hence needs to tediously subtract $1$ until it goes below $2$. This observation leads to the next improvement.

\subsection{Hard-coding the second level}
Per our observation, one obvious improvement would be hard-coding $\alpha_1$ as $\lambda n.(n-2)$. This way its running time will be reduced to $\Theta(1)$, as opposed to the previous $\Theta(n)$.
\begin{lstlisting}
Definition sub_2 (n : nat) : nat
:= match n with | 0 => 0 | 1 => 1 | S (S n') => n' end.
\end{lstlisting}
Without loss of generality, let us assume from now on the constant factors in $\runtime(\alpha_1, n)$ and \cref{lem: inv-ack-hier-runtime} are both $1$. The running time for $\alpha_2$ is then
\begin{equation*}
\runtime\big(\alpha_2, n\big)
 = \sum_{i=0}^{\lceil (n-3)/2 \rceil} \runtime \big(\alpha_1, n-2i\big) + \left\lceil \frac{n-3}{2} \right\rceil =  2\left\lceil \frac{n-3}{2} \right\rceil
 < n 
\end{equation*}
In fact with this improvement, every function in the hierarchy can be computed in linear time, as shown by the next theorem.
\begin{thm} \label{thm: inv-ack-hier-runtime-improved}
	For all $i, n$, $\runtime\big(\alpha_i, n\big) \le 2n + (4^{i+1} - 4)\lceil \log_2n\rceil + 5$.
\end{thm}
Before proving this lemma, we need two crucial technical lemmas.
\begin{lem} \label{lem: inv-ack-3-runtime}
	For all $n$, $\runtime\big(\alpha_3, n\big) \le 2n + 4$.
\end{lem}
\begin{proof}
	It is easy to show that for all $k$, $\alpha_2^{(k)}(n) = \left\lceil \frac{n+3}{2^k} \right\rceil - 3$. Thus
	\begin{equation*}
	\begin{aligned}
	\runtime\big(\alpha_3, n\big) & =
	\sum_{k=0}^{\alpha_{3}(n)}\runtime\left(\alpha_2, \left\lceil \frac{n+3}{2^k} \right\rceil - 3\right) + \alpha_{3}(n) \\
	& \le \sum_{k=0}^{\alpha_3(n)}\left(\frac{n+3}{2^k} + 1\right) - 3\big(\alpha_3(n) + 1\big) + \alpha_3(n) \\
	& \le \sum_{k=0}^{\alpha_3(n)}\frac{n+3}{2^k} - \alpha_3(n) - 2 \le 2(n+3) - 2 = 2n + 4
	\end{aligned}
	\end{equation*}
\end{proof}
\begin{lem}
	For all $i\ge 3$, $\displaystyle \sum_{k=1}^{\alpha_{i+1}(n)} \alpha_i^{(k)}(n) \le 3\big\lceil \log_2n \big\rceil \ \forall n$.
\end{lem}
\begin{proof}
	Let the LHS be $S_i(n)$. Firstly, consider $i = 3$. Note that for $n\le 13$, $S_3(n) = 0$ and for $n\ge 14$, i.e. $\alpha_3(n)\ge 2$, $S_3(n) = \alpha_3(n) + S_3\big(\alpha_3(n)\big)$. The result thus holds for $n\le 13$. Suppose it holds for all $m < n$, where $n\ge 14$. Then
	\begin{equation*}
	S_3(n) \le \alpha_3(n) + 3\big\lceil \log_2(\alpha_3(n)) \big\rceil \le \big\lceil \log_2n \big\rceil + 3\big\lceil \log_2\log_2n \big\rceil
	\end{equation*}
	It is easy to prove $2\big\lceil \log_2\log_2n \big\rceil \le \big\lceil \log_2n \big\rceil$ via induction on $\big\lceil \log_2n \big\rceil$. Thus $S_3(n)\le 3\big\lceil \log_2n \big\rceil$, as desired. Now for $i \ge 4$,
	\begin{equation*}
	S_i(n) = \sum_{k=1}^{\alpha_{i+1}(n)} \alpha_i^{(k)}(n) \le
	\sum_{k=1}^{\alpha_{i+1}(n)} \alpha_3^{(k)}(n) \le
	\sum_{k=1}^{\alpha_{4}(n)} \alpha_3^{(k)}(n) \le 3\big\lceil \log_2n \big\rceil
	\end{equation*}
	The proof is complete.	
%	Let $P(n) \triangleq 2\big\lceil \log_2\log_2n \big\rceil \le \big\lceil \log_2n \big\rceil$. It suffices to prove $P(n) \ \forall n$. Observe that $P(n)$ holds for $n\ge 4$.
\end{proof}
We are ready to present the proof for \Cref{thm: inv-ack-hier-runtime-improved}.
\begin{proof}[Proof of \Cref{thm: inv-ack-hier-runtime-improved}]
	We have proved the result for $i = 0, 1, 2$. Let us proceed with $i\ge 3$ by induction. The base case $i = 3$ has been covered by \cref{lem: inv-ack-3-runtime}. Now suppose the result holds for $i\ge 3$, let $M_i\triangleq 4^{i+1}-4$ for each $i$, we have
	\begin{equation*}
	\begin{aligned}
	& \runtime\big(\alpha_{i+1}, n\big) = \sum_{k=0}^{\alpha_{i+1}(n)} \runtime\left(\alpha_i, \alpha_i^{(k)}(n)\right) + \alpha_{i+1}(n) \\
	& \le \sum_{k=0}^{\alpha_{i+1}(n)}\left(2\alpha_i^{(k)}(n) + M_i\left\lceil\log_2\left(\alpha_i^{(k)}(n)\right)\right\rceil + 5 \right) + \alpha_{i+1}(n) \\
	& \le 2n + M_i\left\lceil\log_2n\right\rceil + 5 + (M_i+2)\sum_{k=1}^{\alpha_{i+1}(n)}\alpha_i^{(k)}(n) + 6\alpha_{i+1}(n) \\
	& \le 2n + M_i\left\lceil\log_2n\right\rceil + 5 +
	3(M_i + 2)\left\lceil\log_2n\right\rceil + 6\left\lceil\log_2n\right\rceil \\
	& = 2n + (4M_i + 12)\left\lceil\log_2n\right\rceil + 5 = 2n + M_{i+1}\left\lceil\log_2n\right\rceil + 5
	\end{aligned}
	\end{equation*}
	, since $4M_i + 12 = 4^{i+2} - 16 + 12 = M_{i+1}$. The proof is complete.
\end{proof}
Thus by hard-coding $\alpha_1$ as $\lambda m.(m-2)$, each level of the inverse Ackermann hierarchy can be computed in $\Theta\big(n + 4^i\log n\big)$ time.

\subsection{An improved inverse Ackermann computation}
Using the previous improvement, we can improve the running time of $\alpha(n)$ per \Cref{defn: inv-ack-hier} by hard-coding the output when $n\le 1 = \Ack(0, 0)$, and starting $\W\alpha$ with $f := \alpha_1$ and $n := \alpha_1(n)$. In other words,
\begin{equation*}
\tilde{\alpha}(n) = \begin{cases}
0 & \text{ if } n \le 1 \\ \W\alpha\big(\alpha_1, \alpha_1(n), 1, n-1\big) & \text{ if } n \ge 2
\end{cases}
\end{equation*}
For $n > 1$, $1\le \min\big\{n-1, \alpha_1(n)\big\}$, so
\begin{equation*}
\W\alpha\big(\alpha_0, \alpha_0(n), 0, n\big) =
\W\alpha\big(\alpha_1, \alpha_1(n), 1, n-1\big)
\end{equation*}
Thus $\tilde{\alpha}(n) = \alpha(n) \ \forall n$. Now, to transform from $\W\alpha\big(\alpha_k, \alpha_k(n), k, n-k\big)$ to $\W\alpha\big(\alpha_{k+1}, \alpha_{k+1}(n), k+1, n-k-1\big)$, $\W\alpha$ needs to compute $\big(\cdt{\alpha_k}{1}\big)\big(\alpha_k(n)\big)$ given $\alpha_k(n)$, which takes time $\runtime\big(\alpha_{k+1}, n\big) - \runtime\big(\alpha_k, n\big)$ by \cref{lem: compose-runtime}; as well as $\alpha_k(n) - k$ for $\Theta(k)$ time. The computation will terminate at $k = \alpha(n)$. Thus for all $n\ge 1$,
\begin{equation*}
\begin{aligned}
& \runtime\big(\tilde{\alpha}, n\big) = \runtime\big(\alpha_1, n\big) + \sum_{k=1}^{\alpha(n) - 1}
\left[ \runtime\big(\alpha_{k+1}, n\big) - \runtime\big(\alpha_k, n\big)
\right] + \sum_{k=1}^{\alpha(n)}\Theta(k) \\
& = \runtime\left(\alpha_{\alpha(n)}, n\right) + \Theta\left(\alpha(n)^2\right)
= \Theta\left(2n + 4^{\alpha(n)}\log_2n + \alpha(n)^2\right) = \Theta(n)
\end{aligned}
\end{equation*}
Therefore, $\tilde{\alpha}$ is able to compute $\alpha$ in linear time.



\section*{Roadmap}
\begin{frame}
\frametitle{Roadmap}
\setlength{\leftmargini}{3.8em}
\setbeamertemplate{enumerate items}[default]
\setbeamercolor{local structure}{fg=black}

\imppar{\textbf{Goal.} Inverting $\Ack$ - without computing $\Ack$.}

\bigskip

\begin{enumerate}[\bfseries Step 1.]
	\itemsep 2ex
	\item<done@1->
	Build the hyperoperations/Ackermann hierarchies via \textbf{Repeater}.
	
	\item<done@1->
	Build inverses of the hyperoperations/Ackermann hierarchies via \textbf{Countdown}, the inverse of \emph{Repeater}.
	
	\item<done@1->
	Use the hierarchy to implement the inverse Ackermann function \\ for inputs encoded in unary. $\bigO (n)$.
	
	\item<come@2-> [\textbf{Bonus.}]
	Improve efficiency for inputs encoded in binary. $\bigO (\log_2n)$.	
\end{enumerate}
\end{frame}


\section{Bonus: The Inverse Ackermann function in logarithmic time for binary encoding system}
\label{sec: binary}
Thus far we have used the Coq type \li{nat}, which represents
a number~$n$ using~$n$ bits.
In contrast, the binary system represents~$n$ in $\lfloor \log_{2} n \rfloor + 1$ bits.
Coq comes with a built-in binary type \li{N}, 
which consists
of constructors \li{N0} and \li{Npos}. The latter, \li{Npos}, unfolds to \li{positive}:

\begin{lstlisting}
Inductive positive : Set :=
  | xI : positive -> positive 
  | xO : positive -> positive
  | xH : positive.
\end{lstlisting}

Constructor \li{xH} represents $1$, and constructors \li{xO} and \li{xI} represent
appending $0$ and $1$ respectively.
By always starting with $1$, \li{positive} bypasses
the issue of disambiguating \emph{e.g.} the strings \li{011} and
\li{00011}, which represent the same number but pose
a minor technical challenge.
To represent $0$, the type \li{N} simply uses the separate constructor \li{N0}.

In both \li{nat} and \li{N}, addition/subtraction of $b$-bit
numbers is $\Theta(b)$, while multiplication is $\Theta \big(b^2\big)$.
In general, arithmetic operations are often faster when the inputs
are encoded in binary. 
In this section we show that this advantage also extends to our techniques.

Our codebase has binary versions of
	\href{https://github.com/inv-ack/inv-ack/blob/7270e64a2600b771f2b1b1b151f7d13fb2ae6c97/bin_repeater.v\#L78-L87}{\color{blue}hyperoperations},
	\href{https://github.com/inv-ack/inv-ack/blob/7270e64a2600b771f2b1b1b151f7d13fb2ae6c97/bin_applications.v\#L30-L36}{\color{blue}inverse hyperoperations},
	\href{https://github.com/inv-ack/inv-ack/blob/7270e64a2600b771f2b1b1b151f7d13fb2ae6c97/bin_repeater.v\#L157-L175}{\color{blue}Ackermann}, and
	\href{https://github.com/inv-ack/inv-ack/blob/7270e64a2600b771f2b1b1b151f7d13fb2ae6c97/bin_inv_ack.v\#L335-L342}{\color{blue}inverse Ackermann}.
Here we show how to compute inverse Ackermann for binary inputs in
$\Theta(b)$ time, where $b$ is the bitlength,
\emph{i.e.} logarithmic time in the input magnitude.
As before, we present an intuitive sketch here and put
full proofs in 
Appendix~D of the extended version of this paper~\cite{extendedinvack}.
% Appendix~\ref{apx:time_analysis_bin}.

\begin{rem}
Although we do not prove it here, our general binary inverse hyperoperations are $O(b^2)$ time,
since Ackermann and base-2 hyperoperations benefit from $\Theta(1)$ division via bitshifts,
whereas general division is $O(b^2)$.
\end{rem}

\subsection{Countdown and contractions in binary}

\renewcommand{\Tleb}{\runtime_{\li{N.leb}}}
\renewcommand{\Tsucc}{\runtime_{\li{N.succ}}}

Although the theoretical \emph{countdown} is independent of the encoding
of its inputs, its Coq definition needs to be adjusted to allow for inputs
in \li{N}. The first step is to translate the arguments of
\li{countdown\_worker} from \li{nat} to \li{N}. Budget \li{b} must
remain in \li{nat} so it can serve as Coq's termination argument,
but all other \li{nat} arguments should be changed
to \li{N}, and functions on \li{nat} to functions on \li{N}.
% Linked by A
%\begin{lstlisting}
%`\href{https://github.com/inv-ack/inv-ack/blob/7270e64a2600b771f2b1b1b151f7d13fb2ae6c97/bin_countdown.v#L104-L109}{Fixpoint bin\_cdn\_wkr}` (f : N -> N) (a n : N) (b : nat) : N :=
%  match b with O => 0 | S b' =>
%    if (n <=? a) then 0 else 1 + bin_cdn_wkr f a (f n) b'
%  end.
%\end{lstlisting}
% Edited by Linh
\begin{lstlisting}
`\href{https://github.com/inv-ack/inv-ack/blob/7270e64a2600b771f2b1b1b151f7d13fb2ae6c97/bin_countdown.v#L104-L109}{\color{blue}Fixpoint bin\_cdn\_wkr}` f a n b : N :=
  match b with O => 0 | S b' =>
    if (n <=? a) then 0
      else 1 + bin_cdn_wkr f a (f n) b'
  end.
\end{lstlisting}

Determining the budget for \li{bin\_countdown\_to} is tricky.
A naïve approach is to use the built-in \li{nat} translation of \li{n},
\emph{i.e.} \li{N.to_nat n}. This is untenable as the translation alone
takes exponential time \emph{viz} the length of $n$'s representation.
We need a linear-time budget calculation for countdowns
of oft-used functions like $\lambda n.(n-2)$.

The key is to focus on functions that can bring their arguments below a threshold via repeated application in
\emph{logarithmic} time, thus allowing a log-sized budget for
\li{bin\_cdn\_wkr}. Simply shrinking by~$1$ is no longer good enough;
we need to halve the argument on every application as shown below:
\begin{defn} \label{defn: bin-contraction}
	$f\in \contract$ is
	\href{https://github.com/inv-ack/inv-ack/blob/7270e64a2600b771f2b1b1b151f7d13fb2ae6c97/bin_countdown.v#L37-L41}{\color{blue}\emph{binary strict above}}
	$a\in \mathbb{N}$ if $\forall n > a, f(n) \le \lfloor \frac{n + a}{2} \rfloor$.
\end{defn}
\noindent The key advantage of binary strict contractions is that if a contraction $f$ is binary strict above some $a$, then \lb
we know that $\forall n > a, \forall k.~f(n) \le \left\lfloor \frac{n - a}{2^k} \right\rfloor + a$.
Therefore, within $\lfloor \log_2 (n - a) \rfloor + 1$ applications of $f$, the result will become equal to or less than~$a$. We can choose this number as the budget for \li{bin\_cdn\_wkr} to successfully reach the countdown value before terminating.
Note that this budget is simply the length of the binary representation
of \li{n - a}, which we calculate using our function
\href{https://github.com/inv-ack/inv-ack/blob/7270e64a2600b771f2b1b1b151f7d13fb2ae6c97/bin_prelims.v#L135-L143}{\color{blue}\li{nat\_size}}. % Linked by A
The Coq definition of countdown on \li{N} is:
\hide{
\begin{lstlisting}
`\href{https://github.com/inv-ack/inv-ack/blob/7270e64a2600b771f2b1b1b151f7d13fb2ae6c97/bin_prelims.v#L135-L143}{\color{blue}Definition nat\_size}` (n : N) : nat :=
  match n with
  | 0 => 0%nat
  | Npos p => let fix nat_pos_size (x : positive) : nat :=
                  match x with
                  | xH => 1%nat
                  | xI y | xO y => S (nat_pos_size y)
                  end
              in nat_pos_size p
  end.
\end{lstlisting}
Note that \li{nat_size} outputs $0$ on $0$, and on any positive number $m$ the size of its binary representation, thus equals to $\lfloor \log_{2} m \rfloor + 1$.
}% end hide
% Linked by A
%\begin{lstlisting}
%`\href{https://github.com/inv-ack/inv-ack/blob/7270e64a2600b771f2b1b1b151f7d13fb2ae6c97/bin_countdown.v#L111-L112}{Definition bin\_countdown\_to}` (f : N -> N) (a n : N) : N :=
%  bin_cdn_wkr f a n (nat_size (n - a)).
%\end{lstlisting}
% Edited by Linh
\begin{lstlisting}
`\href{https://github.com/inv-ack/inv-ack/blob/7270e64a2600b771f2b1b1b151f7d13fb2ae6c97/bin_countdown.v#L111-L112}{\color{blue}Definition bin\_countdown\_to}` f a n :=
  bin_cdn_wkr f a n (nat_size (n - a)).
\end{lstlisting}

\pagebreak
\noindent The following is the binary version of Lemma~\ref{lem: cdt-runtime}:
%\begin{lem} \label{lem: cdt-runtime-bin}
%	$\forall n \in \li{N}$, if $f$ is a binary strict contraction above $a$,
%	\begin{equation*}
%	\runtime_{\cdt{f}{a}}(n) \le \sum_{i=0}^{\cdt{f}{a}(n) - 1} \hspace{-6pt}
%	\runtime_f\big(f^{(i)}(n)\big) \ + \ (\log_2a + 3)\left(\cdt{f}{a}(n) + 1\right) \ + \ 2\log_2n \ + \ \log_2\cdt{f}{a}(n)
%	\end{equation*}
%\end{lem}
% Edited by Linh
\begin{lem} \label{lem: cdt-runtime-bin}
	$\forall n \in \li{N} \ \forall a\in \li{N}$, if $f$ is a binary strict contraction above $a$,
	\begin{equation*}
	\begin{aligned}
	\runtime_{\cdt{f}{a}}(n) \ \le \ \sum_{i=0}^{\cdt{f}{a}(n) - 1} \hspace{-6pt}
	\runtime_f\big(f^{(i)}(n)\big) \ & + \ (\log_2a + 3)\left(\cdt{f}{a}(n) + 1\right) \\ 
  & + \ 2\log_2n \ + \ \log_2\cdt{f}{a}(n)
	\end{aligned}
	\end{equation*}
\end{lem}

%Substituting $a=1$ into \eqref{eq: cdt-runtime-bin} shows that \Cref{lem: inv-ack-hier-runtime} still holds.
%Similar to \Cref{sect: hardcode-lvl2}, the use of binaries is not immediately effective since the first level. We delve deeper into the hierarchy by \emph{hardcoding the $3^{\text{th}}$ level} and starts from there. Now
%\begin{equation*}
%\forall n, n+2  < n+3 < 2(n+2) \iff \forall n,
%\lfloor \log_2(n+2) \rfloor < \lceil \log_2(n+3) \rceil \le \lfloor \log_2(n+2) \rfloor + 1
%\end{equation*}
%This shift from floor to ceiling enables a direct computation, since $\lceil \log_2(n+3) \rceil$ can now be seen as the size of $(n+2)$'s binary representation.
%\begin{lstlisting}
%Definition alpha3 (n: N) : N := N.size (n+2) - 3.
%\end{lstlisting}
%Let $n\ge 2$. The computation of \li{N.size(n)} takes time equal to itself, so the above definition gives $\runtime(\alpha_3, n) \le 2\log_2n$. Fix an $i\ge 3$ and suppose $\runtime(\alpha_i, n) \le C_i\log_2n$. By \Cref{lem: inv-ack-hier-runtime},
\subsection{Inverse Ackermann in $\bigO\left(\log_2 n\right)$}
Our new Coq definition computes countdown
only for strict binary contractions. Fortunately, starting
from $n = 2$, the inverse hyperoperations $a\angle{n}b$ when $a\ge 2$
and the inverse Ackermann hierarchy $\alpha_n$ are all strict binary contractions.
We can construct these hierarchies by hardcoding their
first three levels and recursively building higher levels with \li{bin\_countdown\_to}.
Furthermore,
analagously to the optimization for \li{nat} discussed in~\S\ref{sect: hardcode-lvl2}, we hardcode an additional level.
% Linked by A
%\begin{lstlisting}
%`\href{https://github.com/inv-ack/inv-ack/blob/7270e64a2600b771f2b1b1b151f7d13fb2ae6c97/bin_inv_ack.v#L55-L60}{Fixpoint bin\_alpha}` (m : nat) (x : N) : N :=
%  match m with
%  | 0%nat => x - 1          | 1%nat => x - 2
%  | 2%nat => N.div2 (x - 2) | 3%nat => N.log2 (x + 2) - 2
%  | S m'  => bin_countdown_to (bin_alpha m') 1 (bin_alpha m' x)
%  end.
%\end{lstlisting}
% Edited by Linh
\begin{lstlisting}
`\href{https://github.com/inv-ack/inv-ack/blob/7270e64a2600b771f2b1b1b151f7d13fb2ae6c97/bin_inv_ack.v#L55-L60}{\color{blue}Fixpoint bin\_alpha}` (m : nat) (x : N) : N :=
  match m with
  | 0%nat => x - 1          
  | 1%nat => x - 2
  | 2%nat => N.div2 (x - 2) 
  | 3%nat => N.log2 (x + 2) - 2
  | S m'  => bin_countdown_to
               (bin_alpha m') 1 (bin_alpha m' x)
  end.
\end{lstlisting}
Note that for all $x$, $\li{N.div2}(x - 2) = \left\lfloor \frac{x - 2}{2} \right\rfloor = \left\lceil \frac{x - 3}{2} \right\rceil$ and $\li{N.log2}(x + 2) - 2 = \left\lfloor \log_2(x+2) \right\rfloor - 2 = \left\lceil \log_2(x+3) \right\rceil - 3$, so the above Coq definition is correct.

%\begin{thm} \label{thm: inv-ack-runtime-bin}
%	$\forall i,~\forall n,~\runtime_{\alpha_i}(n) \le 2\log_2n + \left(3\cdot 2^i - 3i - 13\right)\log_2\log_2n + 3i$.
%\end{thm}
% Edited by Linh
\begin{thm} \label{thm: inv-ack-runtime-bin}
	% For all $i$ and $n$,
	\begin{equation*}
	\forall i,n.~\runtime_{\alpha_i}(n) \le 2\log_2n + \left(3\cdot 2^i - 3i - 13\right)\log_2\log_2n + 3i.
	\end{equation*}
\end{thm}

\noindent For any level of the Ackermann hierarchy, this theorem demonstrates
a linear computation time up to the size of the representation of the input, \emph{i.e.} logarithmic time up to its magnitude $n$:
$\runtime_{\alpha_i}(n) = \bigO\big(\log_2n + 2^i\log_2\log_2n \big)$.

Moving on to inverse Ackermann itself, we follow a style nearly identical to that
in~\S\ref{subsec: inv_ack_hier}. For the worker, we simply translate to
\li{N}, keeping the budget in \li{nat} as described earlier.
The inverse Ackermann has an extra hardcoded level.
% Linked by A
%\begin{lstlisting}
%`\href{https://github.com/inv-ack/inv-ack/blob/7270e64a2600b771f2b1b1b151f7d13fb2ae6c97/bin_inv_ack.v#L326-L333}{Fixpoint bin\_inv\_ack\_wkr}` (f : N -> N) (n k : N) (b : nat) : N :=
%  match b with 0%nat  => k | S b' => if n <=? k then k else
%    let g := (bin_countdown_to f 1) in
%      bin_inv_ack_wkr (compose g f) (g n) (N.succ k) b'
%  end.
%
%`\href{https://github.com/inv-ack/inv-ack/blob/7270e64a2600b771f2b1b1b151f7d13fb2ae6c97/bin_inv_ack.v#L335-L342}{Definition bin\_inv\_ack}` (n : N) : N :=
%  if (n <=? 1) then 0 else if (n <=? 3) then 1
%    else if (n <=? 7) then 2 else
%      let f := (fun x => N.log2 (x + 2) - 2) in
%        bin_inv_ack_wkr f (f n) 3 (nat_size n).
%\end{lstlisting}
% Edited by Linh
\begin{lstlisting}
`\href{https://github.com/inv-ack/inv-ack/blob/7270e64a2600b771f2b1b1b151f7d13fb2ae6c97/bin_inv_ack.v#L326-L333}{\color{blue}Fixpoint bin\_inv\_ack\_wkr}` f n k b :=
  match b with 0%nat => k | S b' =>
    if n <=? k then k else
      let g := (bin_countdown_to f 1) in
        bin_inv_ack_wkr
          (compose g f) (g n) (N.succ k) b'
  end.

`\href{https://github.com/inv-ack/inv-ack/blob/7270e64a2600b771f2b1b1b151f7d13fb2ae6c97/bin_inv_ack.v#L335-L342}{\color{blue}Definition bin\_inv\_ack}` n :=
  if (n <=? 1) then 0 else if (n <=? 3) then 1
    else if (n <=? 7) then 2 else
      let f := (fun x => N.log2 (x + 2) - 2) in
        bin_inv_ack_wkr f (f n) 3 (nat_size n).
\end{lstlisting}
% Fixpoint bin_inv_ack_worker (f : N -> N) (n k : N) (b : nat) : N :=
%   match b with
%   | 0%nat  => k
%   | S b' => if n <=? k then k
%             else let g := (bin_countdown f 1) in
%                  bin_inv_ack_worker (compose g f) (g n) (N.succ k) b'
%   end.

% Definition bin_inv_ack (n : N) : N :=
%   if (n <=? 1) then 0
%   else if (n <=? 3) then 1
%        else if (n <=? 7) then 2
%             else let f := (bin_alpha 3) in
%                  bin_inv_ack_worker f (f n) 3 (nat_size n).
\noindent Note that, for $n > 7$, $n < \Ack\big(\lfloor \log_2n \rfloor + 1\big)$ $= \Ack\big(\li{nat_size}(n)\big)$, so a budget of $\li{nat_size}(n)$ suffices.
We show the
\href{https://github.com/inv-ack/inv-ack/blob/7270e64a2600b771f2b1b1b151f7d13fb2ae6c97/bin_inv_ack.v#L437-L472}
{\color{blue}correctness} and
\href{https://github.com/inv-ack/inv-ack/blob/195209ba895061fc51368c3c46c1d8760f05df50/bin_test_runtime_ocaml.ml#L359-L363}
{\color{blue}benchmark}
of \li{bin_inv_ack} in our codebase. Figure~\ref{fig:standalone_binary} 
shows a standalone binary-specific computation of the inverse 
Ackermann function that takes logarithmic 
time up to the input's magnitude.
This is simply an assimilation of the code snippets we have already discussed in 
this section, and serves as a 
binary translation of Figure~\ref{fig:standalone} shown earlier.

\begin{figure}
\lstset{style=myTinyStyle}
\begin{lstlisting}
Require Import Omega Program.Basics.
Open Scope N_scope.

`\href{https://github.com/inv-ack/inv-ack/blob/7270e64a2600b771f2b1b1b151f7d13fb2ae6c97/inv_ack_standalone.v#L56-L64}{\color{blue}Definition nat\_size}` n :=
  match n with
  | 0 => 0%nat
  | Npos p =>
      let fix nat_pos_size (x : positive) : nat :=
        match x with xH => 1%nat
        | xI y | xO y => S (nat_pos_size y) end
      in nat_pos_size p
  end.

`\href{https://github.com/inv-ack/inv-ack/blob/7270e64a2600b771f2b1b1b151f7d13fb2ae6c97/bin_countdown.v#L104-L109}{\color{blue}Fixpoint bin\_cdn\_wkr}` f a n b : N :=
  match b with O => 0 | S b' =>
    if (n <=? a) then 0
      else 1 + bin_cdn_wkr f a (f n) b'
  end.

`\href{https://github.com/inv-ack/inv-ack/blob/7270e64a2600b771f2b1b1b151f7d13fb2ae6c97/inv_ack_standalone.v#L75-L76}{\color{blue}Definition bin\_countdown\_to}` f a n := 
bin_cdn_wkr f a n (nat_size (n - a)).

`\href{https://github.com/inv-ack/inv-ack/blob/7270e64a2600b771f2b1b1b151f7d13fb2ae6c97/inv_ack_standalone.v#L97-L104}{\color{blue}Fixpoint bin\_inv\_ack\_wkr}` f n k b :=
  match b with 0%nat => k | S b' => 
    if n <=? k then k else 
      let g := (bin_countdown_to f 1) in
        bin_inv_ack_wkr
          (compose g f) (g n) (N.succ k) b'
  end.

`\href{https://github.com/inv-ack/inv-ack/blob/7270e64a2600b771f2b1b1b151f7d13fb2ae6c97/inv_ack_standalone.v#L111-L116}{\color{blue}Definition bin\_inv\_ack}` n :=
  if (n <=? 1) then 0 else if (n <=? 3) then 1
    else if (n <=? 7) then 2 else 
      let f := (fun x => N.log2 (x + 2) - 2) in
        bin_inv_ack_wkr f (f n) 3 (nat_size n).
\end{lstlisting}
\caption{A log-time Coq computation of the inverse Ackermann function for inputs represented in binary, \emph{i.e.} \li{N}.}
\label{fig:standalone_binary}
\end{figure}
\lstset{style=myStyle}

As in Theorem~\ref{thm: inv-ack-hardcode-correct}, the time complexity $\runtime_\alpha(n)$ is the sum of each component's runtime:
%\begin{equation*}
%\runtime_\alpha(n) =
%\begin{aligned}
%& \ \runtime_{\alpha_3}(n)
%+ \sum_{k = 3}^{\alpha(n) - 1} \hspace{-5pt} \runtime_{\cdt{\alpha_k}{1}\ }(\alpha_k(n))
%+ \sum_{k = 3}^{\alpha(n)}\Tleb\left(\alpha_k(n), k\right)
%+ \sum_{k = 3}^{\alpha(n) - 1} \hspace{-5pt} \Tsucc(k)\\
%& \ + \ \runtime_{\li{nat_size}}(n)
%\ + \Tleb(n, 1) \ + \Tleb(n, 3) \ + \Tleb(n, 7)
%\end{aligned}
%\end{equation*}
% Edited by Linh
\begin{equation*}
\begin{aligned}
\runtime_{\alpha_3}(n)
& \ + \sum_{k = 3}^{\alpha(n) - 1} \hspace{-5pt} \runtime_{\cdt{\alpha_k}{1}\ }(\alpha_k(n))
+ \sum_{k = 3}^{\alpha(n)}\Tleb\left(\alpha_k(n), k\right) \\
& \ + \sum_{k = 3}^{\alpha(n) - 1} \hspace{-5pt} \Tsucc(k)
\ + \ \runtime_{\li{nat_size}}(n) \\
& \ + \Tleb(n, 1) \ + \ \Tleb(n, 3) \ + \ \Tleb(n, 7)
\end{aligned}
\end{equation*}
With reference to Lemmas~47 (Appendix~C), 52 and 53 (Appendix~D), we have:
% With reference to Lemmas \ref{lem: compose-runtime} (Appendix~C), \ref{lem: leb-runtime-bin} and \ref{lem: succ-runtime-bin} \lb (Appendix~D), we have:
In the second summand, \lb $\runtime_{\cdt{\alpha_k}{1}\ }(\alpha_k(n)) = \runtime_{\alpha_{k+1}}(n) - \runtime_{\alpha_k}(n)$ for each $k$ by Lemma~47.
% \ref{lem: compose-runtime}.
By Lemma~52, 
% By Lemma~\ref{lem: leb-runtime-bin}, 
each $\Tleb$ in the third summand is $\Theta\left(\log_2k\right)$, totalling $\bigO\big(\alpha(n)\log_2\alpha(n)\big) = \text{o}\big(\log_2n\big)$.
The fourth summand is $\Theta(\alpha(n)) = \text{o}\big(\log_2n\big)$ by 
Lemma~53. 
% Lemma~\ref{lem: succ-runtime-bin}. 
The remaining items total $\Theta(\log_2n)$. Thus, $\forall n\ge 8$:
%\begin{equation*}
%\begin{aligned}
%\runtime_\alpha(n)
%& = \runtime_{\alpha_3}(n)
%+ \sum_{k = 3}^{\alpha(n) - 1} \big(\runtime_{\alpha_{k+1}}(n) - \runtime_{\alpha_{k}}(n)\big) + \Theta\left(\log_2 n\right)
%= \runtime_{\alpha_{\alpha(n)}}(n) + \Theta\big(\log_2 n\big) \\
%& = \bigO\big(\log_2n + 2^{\alpha(n)}\log_2\log_2n \big) + \Theta\big(\log_2n\big)
%= \Theta\big(\log_2n\big)
%\end{aligned}
%\end{equation*}
% Edited by Linh
\begin{equation*}
\begin{aligned}
\runtime_\alpha(n)
& = \runtime_{\alpha_3}(n)
+ \sum_{k = 3}^{\alpha(n) - 1} \big(\runtime_{\alpha_{k+1}}(n) - \runtime_{\alpha_{k}}(n)\big) + \Theta\left(\log_2 n\right) \\
& = \runtime_{\alpha_{\alpha(n)}}(n) + \Theta\big(\log_2 n\big) \\
& = \bigO\big(\log_2n + 2^{\alpha(n)}\log_2\log_2n \big) + \Theta\big(\log_2n\big) \\
& = \Theta\big(\log_2n\big)
\end{aligned}
\end{equation*}

\renewcommand{\Tleb}{\runtime_{\li{leb}}}
\renewcommand{\Tsucc}{\runtime_{\li{succ}}}

%We bound the RHS by bounding $\runtime\left(\alpha_{\alpha(n)}(n)\right)$, and $\runtime\left(\alpha_i(n)\right)$ in general. Note that by $\alpha_2(n) = \left\lfloor \frac{n - 2}{2} \right\rfloor$, we have $\runtime(\alpha_2, n) \le 2\log_2 n$ for all $n$.
%Note that this bound can potentially be improved by further tightening the above inequalities. Although we do not obtain an exact asymptotic runtime similar
%to Theorem~\ref{thm: inv-ack-hier-runtime-improved}, since this bound of $4^{\alpha(n)}\log_2n$ is strictly larger than the lower bound of $\log_2n$, it is still extremely small and can be bounded by simpler expressions such as $(\log_2n)^2$.
%Our result is an improved version of the inverse Ackermann function which runs in time $\bigO\big(4^{\alpha(n)}\log_2n\big)$.
%\begin{equation*}
%\alpha(n) \triangleq \begin{cases}
%0 & \text{if } n\le 1\\ 1 & \text{if } 2\le n\le 3 \\ 2 & \text{if } 4\le n\le 7 \\
%\W\alpha\left(\alpha_3,
%\alpha_3(n), 3, n\right) & \text{if } n\ge 8
%\end{cases} \quad \text{ where } \alpha_3 \triangleq \lambda m. \big(\lfloor \log_2(m+2) \rfloor - 2\big)
%\end{equation*} 


\section*{Roadmap}
\begin{frame}
\frametitle{Roadmap}
\setlength{\leftmargini}{3.8em}
\setbeamertemplate{enumerate items}[default]
\setbeamercolor{local structure}{fg=black}

\imppar{\textbf{Goal.} Inverting $\Ack$ - without computing $\Ack$.}

\bigskip

\begin{enumerate}[\bfseries Step 1.]
	\itemsep 2ex
	\item<done@1->
	Build the hyperoperations/Ackermann hierarchies via \textbf{Repeater}.
	
	\item<done@1->
	Build inverses of the hyperoperations/Ackermann hierarchies via \textbf{Countdown}, the inverse of \emph{Repeater}.
	
	\item<done@1->
	Use the hierarchy to implement the inverse Ackermann function \\ for inputs encoded in unary. $\bigO (n)$.
	
	\item<done@1-> [\textbf{Bonus.}]
	Improve efficiency for inputs encoded in binary. $\bigO (\log_2n)$.	
\end{enumerate}
\end{frame}

%\section{Further Discussion}
%\label{sec: discussion}
%\newcommand{\ackt}{\ensuremath{\hat{\alpha}}}

\subsection{The value of a linear-time solution to the hierarchy}

Our functions' linear runtimes can be understood in two distinct but
complementary ways.  A runtime less than the bitlength is impossible
without prior knowledge of the size of the input.  Accordingly, in
an information-theory or pure-mathematical sense, our definitions are
optimal up to constant factors.  And of course in practice, linear-time
solutions are highly usable in real computations.

Sublinear solutions are possible with \emph{a priori} knowledge about
the function and bounds on the inputs one will receive.
An extreme case is $\alpha(n)$, which has value $4$ for all practical
inputs greater than $61$. Accordingly,
this function can be inverted in $O(1)$ in practice.  That said, 
such solutions require external knowledge of the problems and
lookup tables within the algorithm to store precomputed
values, and thus fall more into the realm of engineering than mathematics. 

\subsection{The two-parameter inverse Ackermann function}
Some authors~\cite{chazelle,tarjan} prefer a two-parameter inverse Ackermann function.
\begin{defn} \label{defn: 2para-alpha}
	The two-parameter inverse Ackermann function is defined as:
	\begin{equation} \label{eq: tmp-2para-alpha}
	\ackt (m, n) \triangleq \min\left\{i \ge 1 : \Ack\left(i, \left\lfloor \frac{m}{n} \right\rfloor \right)\ge \log_2n \right\}
	\end{equation}
\end{defn}
Note that $\ackt(n, n)$ and the single-parameter $\alpha(n)$
are neither equal nor directly related, but
it is straightforward to modify our techniques to compute $\ackt(m, n)$.
\hide{This function arises from deep runtime analysis of the disjoint-set data structure. Tarjan \cite{tarjan} showed that, in the disjoint-set data structure, the time required $t(m,n)$ for a sequence of $m$ \textsc{\color{magenta}FIND}s intermixed with $n-1$ \textsc{\color{magenta}UNION}s (such that $m \geq n$) is bounded as: $k_{1}m\cdot\alpha(m,n) \leq t(m,n) \leq k_{2}m\cdot\alpha(m,n)$. In graph theory, Chazelle \cite{chazelle} showed that the minimum spanning tree of a connected graph with $n$ vertices and $m$ edges can be found in time $O(m\cdot\alpha(m,n))$. Computing this function is in fact easier than $\alpha(n)$, as when $m$ and $n$ are given, we are reduced to finding the minimum $i\ge 1$ such that $\Ack_i(s)\ge t$ for $s, t$ fixed, which can be done with the following variant of the \emph{inverse Ackermann worker}.
}
\begin{defn} \label{defn: 2para-inv-ack-worker}
	The {two-parameter inverse Ackermann worker}
	,written $\ackt^{\W}$, is a function $\mathbb{N}^4\to \mathbb{N}$, defined by:
	\hide{$(\mathbb{N}\to \mathbb{N}) \times \mathbb{N}^3\to \mathbb{N}$ such that for all $n, k, b\in \mathbb{N}$ and $f:\mathbb{N}\to \mathbb{N}$:}
%	\begin{equation} \label{eq: 2para-inv-ack-worker-recursion}
%	\ackt^{\W}(f, n, k, b) = \begin{cases}
%	0 & \text{if } b = 0 \vee n\le k \\ 1 + \ackt^{\W}\big(\cdt{f}{1}\circ f , \cdt{f}{1}(n), k, b-1\big) & \text{if } b \ge 1 \wedge n \ge k+1
%	\end{cases}
%	\end{equation}
  \begin{equation} \label{eq: 2para-inv-ack-worker-recursion}
  \begin{aligned}
  & \ackt^{\W}(f, n, k, b) \\
  & \triangleq \begin{cases}
  0 & \text{if } b = 0 \vee n\le k \\ 1 + \ackt^{\W}\big(\cdt{f}{1}\circ f , \cdt{f}{1}(n), k, b-1\big) & \text{if } b \ge 1 \wedge n \ge k+1
  \end{cases}
  \end{aligned}
  \end{equation}
\end{defn}
%Similar to the one-parameter version, the following theorem establishes the correct setting for $\W\alpha_2$ to compute $\alpha(m, n)$.
%\begin{thm}
%	$\displaystyle \ackt(m, n) = 1 + \ackt^{\W}\left(\alpha_1, \alpha_1\big(\lceil\log_2n \rceil\big), \left\lfloor \frac{m}{n} \right\rfloor, \lceil\log_2n \rceil \right)$.
%\end{thm}
% Edited by Linh
\begin{thm} For all $m$ and $n$,
	\begin{equation*}
	\displaystyle \ackt(m, n) = 1 + \ackt^{\W}\left(\alpha_1, \alpha_1\big(\lceil\log_2n \rceil\big), \left\lfloor \frac{m}{n} \right\rfloor, \lceil\log_2n \rceil \right).
	\end{equation*}
\end{thm}
We mechanize the above for both \href{https://github.com/inv-ack/inv-ack/blob/7270e64a2600b771f2b1b1b151f7d13fb2ae6c97/inv_ack.v#L245-L248}{\color{blue}unary} and \href{https://github.com/inv-ack/inv-ack/blob/7270e64a2600b771f2b1b1b151f7d13fb2ae6c97/bin_inv_ack.v#L222-L228}{\color{blue}binary} inputs in our codebase.

\hide{
	\begin{proof}[Proof Sketch]
		It is easy to prove in a similar fashion to \cref{lem: inv-ack-worker-intermediate} that for all $n, b, k$ and $i$, if $\alpha_i(n) > k$ and $b > i$, then
		\begin{equation*}
		\W\alpha_2\big(\alpha_1, \alpha_1(b), k, b\big) = i + \W\alpha_2\big(\alpha_{i+1}, \alpha_{i+1}(n), k, b - i\big)
		\end{equation*}
		Now let $k \triangleq \lfloor m/n \rfloor$, $b \triangleq \lceil \log_2n \rceil$ and $l \triangleq \min\big\{i : \alpha_i(b)\le k\big\}$, which exists because $\Ack(i, \cdot)$ increases strictly with $i$. Then $\alpha(m, n) = \max{1, l}$. If $l = 0$, $\alpha_1(b) \le \alpha_0(b) \le k$, so $\W\alpha_2\big(\alpha_1, \alpha_1(b), k, b\big) = 0$, as desired. If $l \ge 1$,
		\begin{equation*}
		1 + \W\alpha_2\big(\alpha_1, \alpha_1(b), k, b\big)
		= 1 + l - 1 + \W\alpha_2\big(\alpha_l, \alpha_l(b), k, b-l\big) = l
		\end{equation*}
		Here $b\ge l$ due to the fact that $\Ack(b, k)\ge b$, so $\alpha_b(b)\le k$. This completes the proof.
\end{proof}}%end hide




\section*{Conclusion}
\begin{frame}
\frametitle{Conclusion}
%\frametitle{Abstract}
%TODO: Past tense
We inverted the hyperoperation 
and Ackermann hierarchies \\ and implemented 
our inverse hierarchies in Gallina.
We used our inverses to compute the upper inverse of the 
diagonal Ackermann function. 

\bigskip

Our computations run in linear time
-- $\Theta(b)$, where $b$=bitlength -- \\ for inputs represented in 
both unary and binary.

\bigskip

We showed that these functions are consistent with the usual \\ definition of the inverse Ackermann function~$\alpha(n)$.
\keywords{Ackermann  \and hyperoperations \and Coq}
\end{frame}

\section*{Demo}
\begin{frame}
\frametitle{Demo}
We present a brief demonstration of our time bounds in Coq.
\end{frame}

\newcounter{totalframenumber}
\setcounter{totalframenumber}{\inserttotalframenumber}

\appendix


\addtocounter{framenumber}{\value{totalframenumber}}

\setbeamertemplate{footline}{%
	\hspace{1em}\small\arabic{framenumber}/\arabic{totalframenumber}
	\hfill\usebeamertemplate***{navigation symbols}}

\section{Appendix}

\subsection{Runtime of Countdown}
\begin{frame}
\frametitle{Runtime of Countdown}

\textbf{Sum by component in each recursive step.}
\begin{equation*}
\begin{array}{l|llllll}
Step & \multicolumn{2}{l}{Initial \ Arguments} &  & \big(f^{(0)}(n), & n - 0 \big) &  \\[3pt]
\pause 0    & n - 0 > 0,     & f^{(0)}(n) > a & \to & \big(f^{(1)}(n), & n - 1\big) & + 1 \\[3pt]
\pause 1    & n - 1 > 0, & f^{(1)}(n) > a & \to & \big(f^{(2)}(n), & n - 2\big) & + 2 \\[3pt]
\pause \cdots & \cdots & \cdots & \cdots & \cdots & \cdots & \cdots \\
k-1  & n - (k-1) > 0, & f^{(k-1)}(n) > a & \to & \big(f^{(k)}(n), & n - k\big) & + k \\[3pt]
\pause k & n - k \mathrel{?} 0, & {\color{red} f^{(k)}(n) \le a}  &\to & \multicolumn{2}{l}{0} &  \\[3pt] \hline
SUM & \pause \Theta(k) & \pause \Theta((a+1)k) \pause & &  \multicolumn{2}{c}{\sum_{i=0}^{k-1}\runtime_f\big(f^{i}(n)\big)} & \pause \Theta(k)
\end{array}
\end{equation*}

%TODO: Increase vertical space after \hline

\pause
Substitute $k = \cdt{f}{a}(n)$:
%\begin{equation*}
\imppar{$\displaystyle
\runtime_{\cdt{f}{a}}(n) = \sum_{i=0}^{\cdt{f}{a}(n)-1}\runtime_f\big(f^{i}(n)\big)
+ \Theta\left((a+1)\cdt{f}{a}(n)\right).
$}
%\end{equation*}
\end{frame}


\subsection{Time Complexity for Unary Inputs: $\bigO(n)$}

\begin{frame}
\frametitle{Runtime of Each $\alpha_i$ (Unary Inputs, Asymptotic Bounds)}

$\displaystyle
\alpha_{i+1} = \cdt{\big(\alpha_i\big)}{1}\circ \alpha_i
\quad \implies \quad \runtime_{\alpha_{i+1}}(n) = \runtime_{\alpha_i}(n) + \runtime_{\cdt{\alpha_i}{1}}\big(\alpha_i(n)\big)
$

\smallskip
\pause 
\imppar{
$\displaystyle
\runtime_{\alpha_{i+1}}(n) = \runtime_{\alpha_i}(n) + 
\sum_{i=0}^{\cdt{\alpha_i}{1}(\alpha_i(n)) - 1} \runtime_{\alpha_i}\left(\alpha_i^{(i+1)}(n)\right) + \Theta\left(\cdt{\alpha_i}{1}(\alpha_i(n)) \right)
$
}

\bigskip

\pause 
Substitute $\cdt{\alpha_i}{1}(\alpha_i(n))$ for $\alpha_{i+1}(n)$:
\bigskip

$ \displaystyle 
\runtime_{\alpha_{i+1}}(n)
  = \runtime_{\alpha_i}(n) + 
  \sum_{i=0}^{\alpha_{i+1}(n) - 1} \runtime_{\alpha_i}\left(\alpha_i^{(i+1)}(n)\right) + \Theta\left(\alpha_{i+1}(n) \right)
$

\pause 
\smallskip
\imppar{
$\displaystyle 
\runtime_{\alpha_{i+1}}(n)
	= \sum_{i=0}^{\alpha_{i+1}(n)} \runtime_{\alpha_i}\left(\alpha_i^{(i)}(n)\right) + \Theta\left(\alpha_{i+1}(n) \right)
$
}

\end{frame}


\begin{frame}
\frametitle{Runtime of Each $\alpha_i$ (Unary Inputs, Precise Bounds)}

Countdown runtime:
\begin{equation*}
\textstyle \runtime_{\cdt{f}{a}}(n) =
\sum_{i=0}^{\cdt{f}{a}(n) - 1} \runtime_f\left(f^{(i)}(n)\right)
+ (a + 2)\cdt{f}{a}(n) + f^{\left(\cdt{f}{a}(n)\right)}(n) + 1
\end{equation*}

\smallskip

\pause 
$\alpha_2$ and $\alpha_3$ runtime: $\runtime_{\alpha_2}(n) \le 2n - 2$ and $\runtime_{\alpha_3}(n) \le 4n + 4$.

\bigskip 

\pause 
$\alpha_i$ runtime:
\begin{equation*}
\textstyle \runtime_{\alpha_{i+1}}(n) \le \sum_{k=0}^{\alpha_{i+1}(n)}\runtime_{\alpha_i}\left( \alpha_i^{(k)}(n)\right) + 3\alpha_{i+1}(n) + 3
\end{equation*}

\smallskip

\pause 
\imppar{
\begin{minipage}{\linewidth}
	\textbf{Theorem.} $\forall i.~\runtime_{\alpha_i}(n) \le 4n + \left(19\cdot 2^{i-3} - 2i - 13\right)\log_2n + 2i = \bigO(n)$, when using $\alpha_1 \triangleq \lambda n. (n-2)$.
\end{minipage}}

\end{frame}



\begin{frame}
\frametitle{Runtime of Inverse Ackermann for Unary Inputs}

\begin{equation*}
\begin{array}{l|llllllll}
Step & \multicolumn{2}{l}{Initial \ Arguments} &  &  \big(\alpha_1, & \alpha_1(n), & 1, & b - 1 \big) \\[3pt]
\pause 1 & b - 1 > 0, & \alpha_1(n) > 1 & \to & \big(\alpha_2, & \alpha_2(n), & 2, & b - 2\big) \\[3pt]
2 & b - 2 > 0, & \alpha_2(n) > 2 & \to  & \big(\alpha_3, & \alpha_3(n), & 3, & b - 3\big) \\
\cdots & \cdots &\cdots & \cdots & \cdots & \cdots & \cdots \\
%b > k-1, \ \alpha_{k-1}(n) > k - 1 & \to  & \big(\alpha_k, & \alpha_k(n), & k, & b - k\big) & k-1 \\
k & b - k > 0, & {\color{red} \alpha_{k}(n) \le k} & \to  & \multicolumn{4}{l}{k = \alpha(n)} \\[3pt] \hline
\pause SUM & \pause \Theta(k) & \pause \Theta\big(\sum_{i=1}^{k}i\big) \pause & & \multicolumn{4}{l}{\sum_{i=1}^{k-1} \runtime_{\cdt{\alpha_i}{1}} (\alpha_i(n))} \\[3pt]
\pause = & \Theta(k) & \Theta(k^2) & & \multicolumn{4}{l}{\sum_{i=1}^{k-1} \big(\runtime_{\alpha_{i+1}}(n) - \runtime_{\alpha_i}(n) \big)}
\end{array}
\end{equation*}

\smallskip

\pause
Therefore, 
\imppar{$\displaystyle
\begin{aligned}
\runtime_{\alpha}(n)
& = \runtime_{\alpha_{\alpha(n)}}(n) - \runtime_{\alpha_1}(n) + \Theta\left(\alpha(n)^2\right) + \runtime_{\alpha_1}(n) \\
& = \bigO \left(n + 2^{\alpha(n)}\log_2n + \alpha(n)^2 \right) = \bigO (n)
\end{aligned}
$}

\end{frame}

\subsection{Time Complexity for Binary Inputs: $\bigO (\log_2n)$}


\begin{frame}
\frametitle{Runtime of Each $\alpha_i$ (Binary Inputs, Precise Bounds)}

Countdown runtime:
\begin{equation*}
\runtime_{\cdt{f}{a}}(n) \le \sum_{i=0}^{\cdt{f}{a}(n) - 1} \hspace{-6pt}
\runtime_f\big(f^{(i)}(n)\big) \ + \ (\log_2a + 3)\left(\cdt{f}{a}(n) + 1\right) \ + \ 2\log_2n \ + \ \log_2\cdt{f}{a}(n)
\end{equation*}

\pause 
$\alpha_3$ runtime: $\runtime_{\alpha_3}(n) \le 2\log_2n + \log_2\log_2n + 3$.

\bigskip

\pause 
$\alpha_i$ runtime:
$
%\begin{equation*}
\runtime_{\alpha_{i+1}}(n) \le
\sum_{k=0}^{\alpha_{i+1}(n)}\runtime_{\alpha_i}\big(\alpha_i^{(k)}(n)\big)
+ 6\log_2\log_2n + 3.
%\end{equation*}
$

\bigskip

\pause 
\textbf{Theorem.} $\forall i. \runtime_{\alpha_i}(n) \le 2\log_2n + \left(3\cdot 2^i - 3i - 13\right)\log_2\log_2n + 3i = \bigO(\log_2n)$, when hard-coding up to $\alpha_3$.

\end{frame}



\begin{frame}
\frametitle{Runtime of Inverse Ackermann for Binary Inputs}

\begin{equation*}
\begin{array}{l|llllllll}
Step & \multicolumn{2}{l}{Initial \ Arguments} &  &  \big(\alpha_3, & \alpha_3(n), & 3, & b - 0 \big) \\[3pt]
\pause 1 & b - 0 > 0, & \alpha_3(n) > 3 & \to & \big(\alpha_4, & \alpha_4(n), & 4, & b - 1\big) \\[3pt]
2 & b - 1 > 0, & \alpha_4(n) > 4 & \to  & \big(\alpha_5, & \alpha_5(n), & 5, & b - 2\big) \\
\cdots & \cdots &\cdots & \cdots & \cdots & \cdots & \cdots \\
%b > k-1, \ \alpha_{k-1}(n) > k - 1 & \to  & \big(\alpha_k, & \alpha_k(n), & k, & b - k\big) & k-1 \\
k & b - k > 0, & {\color{red} \alpha_{k}(n) \le k} & \to  & \multicolumn{4}{l}{k = \alpha(n)} \\[3pt] \hline
SUM & \pause \Theta(k) & \pause \Theta\big(\sum_{i=1}^{k}\log_2i\big) \pause & & \multicolumn{3}{r}{\sum_{i=1}^{k-1} \runtime_{\cdt{\alpha_i}{1}} (\alpha_i(n))} \pause & \Theta(k) \\[3pt]
\pause = & \Theta(k) & \Theta(k\log_2k) & \multicolumn{4}{r}{\sum_{i=1}^{k-1} \big(\runtime_{\alpha_{i+1}}(n) - \runtime_{\alpha_i}(n) \big)} & \Theta(k)
\end{array}
\end{equation*}

\smallskip

\pause 
\imppar{
$\displaystyle
\begin{aligned}
\runtime_{\alpha}(n)
& = \runtime_{\alpha_{\alpha(n)}}(n) - \runtime_{\alpha_3}(n) + \Theta\left(\alpha(n)\log_2 \alpha(n)\right) + \runtime_{\alpha_3}(n) \\
\pause & = \bigO \left(\log_2n + 2^{\alpha(n)}\log_2\log_2n + \alpha(n)\log_2\alpha(n) \right) = \bigO (\log_2n)
\end{aligned}
$}

\end{frame}



%\begin{frame}[allowframebreaks]
%\frametitle{References}
%\bibliographystyle{amsalpha}
%\bibliography{inv_ack}
%\end{frame}

\end{document}